% Zeit und Raum. Grenze, Richtung, Maß
% Master-Datei: bindet die Kapitel aus kapitel/ und den Apparat aus apparat/ ein.
\documentclass[fontsize=11pt,paper=a4,oneside,headings=small,open=any,numbers=noenddot,toc=listof]{scrbook}
\usepackage[T1]{fontenc}
\usepackage[utf8]{inputenc}
\usepackage[ngerman]{babel}
\usepackage{lmodern}
\usepackage{microtype}
\usepackage[autostyle=true,german=guillemets]{csquotes}
\usepackage{amsmath,amssymb}
\usepackage{booktabs}
\usepackage{enumitem}
\usepackage{tikz}
\usetikzlibrary{arrows.meta,calc,decorations.pathmorphing,positioning,patterns,shapes.geometric}
\usepackage[hidelinks]{hyperref}

\setlist{nosep,leftmargin=*}
\setcounter{tocdepth}{0}
\setkomafont{disposition}{\normalfont\bfseries}
\setkomafont{caption}{\small}
\setcapindent{0pt}
\KOMAoptions{captions=tableheading}
\renewcommand*{\chapterheadstartvskip}{\vspace*{1.5\baselineskip}}
\renewcommand*{\chapterheadendvskip}{\vspace{1.2\baselineskip}}

\newcommand{\zitat}[1]{\enquote{#1}}

\title{Zeit und Raum}
\subtitle{Grenze, Richtung, Maß}
\author{}
\date{}

\begin{document}
\frontmatter
\maketitle
\tableofcontents
\mainmatter

% --- Kapitel -----------------------------------------------------------
\chapter{Das Und}

Wer \zitat{Zeit und Raum} sagt, hat schon entschieden, ohne zu wissen wie. Die Konjunktion stellt zwei Namen nebeneinander, als wären sie Nachbarn. Aber Nachbarschaft ist eine räumliche Beziehung. Sagt man stattdessen, Zeit und Raum seien \emph{zugleich} da, so hat man die Zeit zur Bühne gemacht, auf der beide stehen. Eugen Fink hat diese Verlegenheit als erste Frage einer Untersuchung von Raum, Zeit und Bewegung festgehalten: Wie sind die beiden beisammen? Nebeneinander oder gleichzeitig? Keine der beiden Antworten ist neutral, jede nimmt Partei für eine der beiden Seiten, und zwar bevor irgendetwas über die Sache gesagt ist.\footnote{Fink, Zur ontologischen Frühgeschichte von Raum -- Zeit -- Bewegung, Vorverständigung.}

Julius Baumann hat die Parteinahme, die in der Geschichte des Denkens tatsächlich stattgefunden hat, in einem Satz zusammengefasst: \zitat{Die Zeit hat das Schicksal gehabt, gewöhnlich nach dem Raum und als diesem parallel behandelt zu werden.}\footnote{Baumann, Die Lehren von Raum, Zeit und Mathematik, Bd.\,2, Schluss.} Das Wort \zitat{Schicksal} ist genau. Es war keine Entscheidung, die jemand getroffen hätte, sondern eine Gewohnheit, die sich aus der Sache selbst zu ergeben schien: Der Raum liegt vor, er lässt sich zeichnen, messen, teilen und wieder zusammensetzen. Die Zeit liegt nicht vor. Was von ihr da ist, ist ein Jetzt, das schon nicht mehr ist, wenn man es zeigt. Also nimmt man das Vorliegende als Muster für das Nichtvorliegende: eine Linie für die Zeit, Punkte für die Augenblicke, Strecken für die Dauern. Und weil das Muster so gut passt, bemerkt man nicht, was es verdeckt.

Es verdeckt drei Dinge. Es verdeckt, dass die Zeit ihre Grenze anders hat als der Raum: Der Raum ist als Ganzes da und lässt sich überall begrenzen, die Zeit ist nur als Grenze da und lässt sich nirgends beliebig begrenzen. Es verdeckt, dass die Zeit eine Richtung hat, die kein Raum aus sich hervorbringt; wo der Raum Richtungen zeigt, oben und unten, vorn und hinten, rechts und links, stammen sie von einem Leib, der in ihm lebt, und der Leib ist ein Zeitliches. Und es verdeckt, dass das Maß, mit dem der Raum gemessen wird, aus der Zeit kommt: Das Lichtjahr ist keine astronomische Bequemlichkeit, sondern die Form aller Längenmessung, sobald man fragt, was Länge physikalisch bedeutet.

Grenze, Richtung, Maß: Das sind die drei Asymmetrien, um die es im Folgenden geht. Sie werden nicht behauptet, sondern an den Texten aufgesucht, die von der Sache handeln. Diese Texte gehören zwei Lagern an. Das eine Lager ist das der Parallelisierung. Ihm gehören Zenon an, wie Fink ihn liest, Kant, der beide Anschauungsformen in einem Kapitel behandelt, Minkowski, der die Zeit zur weiteren Koordinate macht, und Hegel, der Raum und Zeit \zitat{zwei Weisen des Außersichseins} nennt. Das andere Lager ist das der Asymmetrie: Reichenbach, für den die Zeit die Faserrichtung der Kausalketten ist und der Raum nur ihr Nebeneinander; Brentano, für den das räumliche Kontinuum als Ganzes besteht und das zeitliche nur als Grenze; Guzzoni, für die die Zeit zum Raum wird und nicht umgekehrt; Koch, für den kein Raumpunkt von seinem Spiegelbild unterscheidbar ist, solange nicht ein Leib das Hier besetzt. Und Hegel gehört auch dem zweiten Lager an, denn bei ihm geht der Raum in die Zeit über, und \zitat{die Zeit ist also erst der wahrhafte Punkt}.\footnote{Hegel, Vorlesung über Naturphilosophie, Nachschrift Uexküll, zitiert nach Bonsiepen, Hegels Raum-Zeit-Lehre.}

Dass Hegel in beiden Lagern steht, ist kein Zufall und keine Inkonsequenz. Es zeigt, was die Parallelisierung eigentlich ist: nicht ein Irrtum, sondern eine Abstraktion, die ihre eigene Aufhebung schon enthält. Wer Zeit und Raum parallel behandelt, muss an irgendeiner Stelle sagen, wodurch sie sich unterscheiden. Er legt diese Stelle nur dorthin, wo sie am wenigsten auffällt: in ein Vorzeichen, in eine Fußnote, in einen Übergang. Die Aufgabe ist deshalb nicht, die Parallelisierung zu widerlegen, sondern die Stellen aufzusuchen, an denen sie die Asymmetrie versteckt, und sie dort herauszuholen.

Was dabei herauskommt, ist keine Umkehrung des Schicksals. Es wäre leicht, nun den Raum nach der Zeit und als dieser parallel zu behandeln, und einige der besten Gedanken, die im Folgenden vorkommen, sind in Gefahr, genau das zu tun: Reichenbachs Reduktion der Raummessung auf Zeitmessung, Hegels Übergang des Raumes in die Zeit, Heideggers frühe Rückführung der Räumlichkeit auf die Zeitlichkeit. Aber jeder dieser Gedanken trägt seinen eigenen Vorbehalt bei sich. Reichenbach braucht einen materiellen Messkörper, den die Lichtgeometrie nicht liefert. Kant braucht ein Beharrliches im Raum, ohne das die Zeit nicht einmal als Größe vorgestellt werden könnte. Heidegger hat seine Rückführung zurückgenommen: Sie \zitat{läßt sich nicht halten}.\footnote{Heidegger, Zeit und Sein, zitiert nach Guzzoni, Weile und Weite.} Der Raum widersteht. Was er ist, wenn er nicht das Parallele der Zeit und nicht ihr Abgeleitetes ist, bleibt am Ende offen, und es wird nicht so getan, als wäre es geschlossen.

\chapter{Die Parallelisierung}

\section{Zenon}

Fink liest die Paradoxien Zenons nicht als Sophismen, sondern als die erste systematische Dialektik von Raum, Zeit und Bewegung. Ihre Wurzel ist ein ungeklärtes \zitat{In-sein}. Dass ein Ding im Raum ist, versteht Zenon nach dem Muster des Gefäßes: Es besetzt einen Ort und keinen anderen. Dass ein Ding in der Zeit ist, versteht er nach demselben Muster: Es besetzt ein Jetzt und kein anderes. Fink nennt das ein \zitat{Doppelgeleise}. Auf beiden Gleisen läuft dieselbe Logik, und auf beiden entsteht dieselbe Schwierigkeit: Das Kontinuum ist vor seinen Teilen, aber das Muster des Gefäßes verlangt, dass es aus ihnen bestehe. Der fliegende Pfeil ruht, weil er in jedem Jetzt an einem Ort ist; Achilles erreicht die Schildkröte nicht, weil zwischen zwei Orten unendlich viele liegen.\footnote{Fink, Frühgeschichte, Kapitel zu Zenon.}

Die stillschweigende Voraussetzung, sagt Fink, ist, ein Ding könne nicht zugleich an zwei Orten und nicht am selben Ort zu zwei Zeiten sein. \zitat{In dieser massiven Formulierung ist das sicher richtig, aber Bewegung ist die seltsame Weise, wo ein Ding im selben Moment an zwei Orten ist, d.\,h. im Übergang zwischen zwei Orten sich befindet.}\footnote{Fink, Frühgeschichte, Kapitel zu Zenon.} Was hier auffällt, ist, dass der Übergang zeitlich beschrieben wird (\zitat{im selben Moment}) und räumlich gescheitert ist (\zitat{an zwei Orten}). Die Parallelisierung bricht genau da, wo sie sich bewähren sollte, in der Bewegung, und sie bricht auf der Seite des Raumes. Vom Raum her lässt sich der Übergang nicht denken, weil ein Ort ein Ort ist. Von der Zeit her lässt er sich denken, weil ein Jetzt nie nur ein Jetzt ist; was das heißt, wird sich an der Grenze zeigen, die das Jetzt ist. Das ist die erste Spur der Asymmetrie: Sie liegt nicht in einem Merkmal, das die Zeit hat und der Raum nicht, sondern darin, dass die Zeit den Übergang trägt, den der Raum zerlegt.

\section{Kant}

Reichenbach hat Kant vorgeworfen, Raum und Zeit \zitat{in einem gemeinsamen Kapitel} behandelt zu haben, so als wären es zwei Fälle desselben, die Anschauungsform.\footnote{Reichenbach, Philosophie der Raum-Zeit-Lehre, \S\,16.} Der Vorwurf trifft die transzendentale Ästhetik. Er trifft nicht die Analytik. Marko Martić hat gezeigt, dass Kant in der Analytik einen \emph{synthetischen} Vorrang der Zeit ansetzt: Die Schemata sind Zeitbestimmungen, die Reihe ist zeitlich, der Raum dagegen ist \zitat{ein Aggregat, aber keine Reihe}.\footnote{Kant, Kritik der reinen Vernunft, B\,439, nach Martić, Zeit und Raum.} Was aufgezählt werden soll, muss nacheinander aufgezählt werden, und das Nacheinander ist Zeit. Zugleich hält Kant an einer Gegenrichtung fest, die in der Widerlegung des Idealismus ausgesprochen wird: Die Bestimmung meiner Existenz in der Zeit setzt ein Beharrliches voraus, und das Beharrliche ist nicht in mir, sondern im Raum. In einer Reflexion heißt es, ohne Raum würde die Zeit selbst nicht als Größe vorgestellt werden. Martić folgt Reinhard Brandt in der Folgerung, \zitat{daß der Raum also die Zeit voraussetzt und nur als Zeitraum konzipierbar ist}, und hält zugleich die Umkehrung fest: Die Zeit wird nur am Raum zur Größe.\footnote{Martić, Zeit und Raum, Kapitel zum Schematismus und zur Widerlegung des Idealismus.} Das ist nicht mehr Parallelisierung. Es ist ein Wechselverhältnis, in dem die Zeit die Ordnung stiftet und der Raum das Maß ermöglicht. Es wird sich zeigen, dass Reichenbach dieselbe Struktur aus der Physik gewinnt.

\section{Minkowski}

Die vollständigste Gestalt der Parallelisierung ist die, die die Zeit zur weiteren Koordinate macht. Bei Minkowski wird sie als imaginäre Koordinate eingeführt, damit das Linienelement die Form einer euklidischen Länge annimmt. Ilse Schneider hat schon 1921 festgehalten, was diese Wendung bedeutet und was nicht: Raum und Zeit \zitat{verlieren die physikalische Gegenständlichkeit auch vereint, und nur durch die Materie kommt die physikalische Bedeutung hinein}.\footnote{Schneider, Das Raum-Zeit-Problem bei Kant und Einstein, Kapitel zur allgemeinen Relativitätstheorie.} Die Koordinaten sind Symbole, messbar ist das Linienelement. Damit ist die Zeit zwar der Form nach den Raumkoordinaten gleichgestellt, aber die Gleichstellung ist eine Gleichstellung von Symbolen, nicht von Sachen.

Reichenbach geht einen Schritt weiter. Für ihn ändert die Auffassung der Zeit als weitere Dimension nichts am Sondercharakter der Zeit. Das negative Vorzeichen, das in der metrischen Fundamentalform vor dem Zeitglied steht, ist nicht der Preis, den man für die Gleichstellung zahlt, sondern der mathematische Ausdruck dafür, dass sie nicht gelingt. Und er fügt hinzu, dass das Wesen der Zeit mit dieser Charakterisierung \zitat{noch nicht erschöpft} ist: Das Vorzeichen spricht nur eine Trennung von Raum und Zeit aus, nicht den eigentümlichen Charakter der Zeit.\footnote{Reichenbach, Raum-Zeit-Lehre, \S\,16 und \S\,30.} Ob man das Vorzeichen als solches schreibt oder es in eine imaginäre Einheit versteckt, ist Konvention. Dass eine der Koordinaten anders ist als die anderen, ist keine. An dieser Stelle zeigt sich am deutlichsten, was oben behauptet wurde: Die Parallelisierung legt die Asymmetrie an einen Ort, wo sie nicht auffällt. Sie schafft sie nicht ab.

\section{Hegel}

Hegel nennt Raum und Zeit \zitat{zwei Weisen des Außersichseins}: Der Raum ist das gleichgültige Außersichsein als ruhendes, die Zeit dasselbe in seiner Negativität, als Werden.\footnote{Hegel, Naturphilosophie, Nachschriften, nach Bonsiepen.} Das klingt nach Parallelität, und Bonsiepen hat gezeigt, dass Hegel den Raum tatsächlich in einer Weise denkt, die dem Sein des Parmenides näher steht als dem Werden: als Sichselbstgleichheit, als daseienden Äther. Aber die beiden Weisen sind nicht gleichrangig. Der Punkt, der im Raum die Negation des Raumes sein soll, hat im Raum keine Wirklichkeit; er wird sogleich zur Linie, die Linie zur Fläche. Erst in der Zeit hat die Negativität ihre Wirklichkeit. Deshalb der Satz aus der Nachschrift: \zitat{Die Zeit ist also erst der wahrhafte Punkt.} Und deshalb geht bei Hegel der Raum in die Zeit über, nicht die Zeit in den Raum. Die Parallelisierung ist bei ihm ein Anfang, der sich selbst aufhebt.

Das ist die Lage. Alle, die Zeit und Raum parallel behandelt haben, wissen um den Unterschied. Zenon scheitert an ihm, Kant baut ihn in die Analytik ein, Minkowski versteckt ihn im Vorzeichen, Hegel lässt ihn im Übergang hervortreten. Er wird nun an den drei Stellen herausgeholt, an denen er sich am wenigsten verbergen lässt: an der Grenze, an der Richtung, am Maß.

\chapter{Grenze}

\section{Das Ganze und die Grenze}

Brentano unterscheidet zwei Arten des Kontinuierlichen nach der Weise, wie sie bestehen. Das topisch Kontinuierliche, der Körper, besteht dem Ganzen nach: Alle seine Teile sind zugleich da. Das chronisch Kontinuierliche, das Verlaufende, besteht nur einer Grenze nach. Von einer Melodie ist immer nur ein Ton wirklich, von einer Bewegung nur eine Lage. Was vorher war, ist nicht mehr, was nachher kommt, ist noch nicht. Das Jetzt ist also Grenze eines Kontinuums, das im Übrigen nicht besteht.\footnote{Brentano, Philosophische Untersuchungen zu Raum, Zeit und Kontinuum, Erster Teil, Kapitel zum Kontinuierlichen.}

Man könnte einwenden, eine Grenze ohne das Begrenzte sei nichts, und das Jetzt sei deshalb ebenso nichts wie ein Punkt ohne Linie. Brentano nimmt den Einwand auf und wendet ihn. Die innere Wahrnehmung zeigt in recto die Grenze und in obliquo die Strecke, für die der Punkt eine Grenze ist. \zitat{Fragt man, in welcher Länge die Strecke in obliquo vorgestellt und anerkannt werde, so ist zu antworten: in gar keiner bestimmten Länge.}\footnote{Brentano, Untersuchungen, S.\,131.} Das Jetzt ist also nicht die Grenze einer Strecke, die daneben vorhanden wäre, sondern die Grenze einer Strecke, die nur in der Grenze mitgemeint ist, und zwar ohne Maß. Der Raum verhält sich gerade umgekehrt. Er wird als Ganzes vorgestellt, und seine Grenzen, Flächen, Linien, Punkte, werden an ihm gesetzt. Hegel sagt es in der Nachschrift kurz: \zitat{Man kann einem Raume überall Grenzen setzen, den Raum aber nicht unterbrechen.}\footnote{Hegel, Naturphilosophie, Nachschrift, nach Bonsiepen.} Überall, das heißt: an beliebiger Stelle; nicht unterbrechen, das heißt: keine der gesetzten Grenzen ist eine Grenze des Raumes selbst.

Daraus folgt eine Bestimmung des Körpers, die Brentano ausdrücklich zieht: Ein Körper, der in der Zeit besteht, ist die dreidimensionale Grenze eines Kontinuums, das eine Dimension mehr hat, der Zeit. Er besteht nur als diese Grenze. \zitat{Wenn ein Ding, das besteht, beharrt, so geschieht dies durch ein kontinuierlich sich wiederholendes Sich-Erneuern.}\footnote{Brentano, Untersuchungen, Kapitel zum Beharren.} Das Beharren ist also keine Ruhe, sondern eine Tätigkeit ohne Unterschied. Was räumlich als ruhende Substanz erscheint, ist zeitlich ein Erneuern, das sich von Grenze zu Grenze fortsetzt. Hier zeigt sich zum ersten Mal, dass das, was am Raum am festesten scheint, das Beharrliche, seine Weise zu sein von der Zeit hat.

\section{Der wahrhafte Punkt}

Hegel gelangt von der anderen Seite zum selben Ergebnis. Der Punkt soll im Raum die Negation des Raumes sein, das Fürsichsein, der im Raum seiende Gedanke, wie Bonsiepen paraphrasiert. Aber im Raum hat der Punkt keine Wirklichkeit. Sobald er gesetzt ist, ist er Linie, die Linie ist Fläche, die Fläche ist Raum. Die Negation, die der Punkt sein sollte, wird sogleich wieder ausgebreitet. Deshalb heißt es in der Nachschrift: \zitat{Die Zeit ist also erst der wahrhafte Punkt.} In der Zeit hat die Negativität Wirklichkeit: Das Jetzt ist, indem es nicht ist, und darin ist es, was der Punkt im Raum nur sein soll. Was der Punkt im Raum nicht sein kann, ist er als Jetzt; und ein Ort, so lässt sich der Übergang lesen, ist nicht ein Raumteil, sondern ein Punkt, der als Jetzt gesetzt ist, also die Zeit im Raum.\footnote{Koch, Hegelsche Subjekte in Raum und Zeit, zum Übergang des Raumes in die Zeit; Bonsiepen, Hegels Raum-Zeit-Lehre, zu den Nachschriften. Die Bestimmung des Ortes als gesetztes Jetzt ist eigene Lesart des Übergangs, kein Wortlaut Hegels.}

Fink sagt es ohne die Dialektik: Der Raum ist primär Kontinuum, die Zeit ist als Jetzt primär diskret, \zitat{das Diskrete schlechthin}.\footnote{Fink, Frühgeschichte, Kapitel zur aristotelischen Zeit-Analytik.} Man kann die Ausdehnung des Raumes teilen, aber die Teilung kommt zum Raum hinzu; das Jetzt kommt zur Zeit nicht hinzu, es ist die Weise, wie sie ist.

\section{Ein Widerspruch, der keiner ist}

Nun sagt Baumann das Gegenteil. Für ihn ist gerade das Diskrete, das Abgesetzte in der Zeit \emph{nicht}; die Zeit ist \zitat{nicht aus Momenten als aus Theilen zusammengesetzt}, und die Zahl gehört mehr dem Raum als der Zeit.\footnote{Baumann, Lehren, Bd.\,2, S.\,658--668.} Guzzoni bestätigt ihn von einer ganz anderen Seite: Es gibt keine halbe und keine doppelte Weile. Eine Weile ist eine Phase der Zeit, die nicht als Abfolge von Jetztpunkten erfahren wird, und ihre Dauer ist nur qualitativ bestimmt.\footnote{Guzzoni, Weile und Weite, Kapitel zur Weile.} Wer die Zeit teilt, wie man den Raum teilt, in Abschnitte gleicher Länge, hat sie schon nach dem Raum behandelt.

Fink und Baumann widersprechen einander nur, wenn man \zitat{diskret} in beiden Sätzen gleich versteht. Sie meinen aber zwei Dinge. Fink meint den einen Schnitt, der die Zeit ist: das Jetzt, vor dem alles vergangen und nach dem alles künftig ist. Baumann meint die vielen Schnitte, die man an einer Linie anbringen kann, um sie in Teile zu zerlegen. Die Zeit hat den einen und nicht die vielen. Der Raum hat die vielen und nicht den einen. Brentano hat diese Kreuzung ausgesprochen, ohne sie als solche zu benennen: \zitat{Bei der Zeit [\ldots] ist dieser Punkt der Sache nach ein ausgezeichneter. Beim Raum gibt es keinen Punkt, der in ähnlicher Weise sich der Sache nach unterschiede.}\footnote{Brentano, Untersuchungen, S.\,117.}

\begin{figure}[htbp]
\centering
\begin{tikzpicture}[scale=1]
% --- Raum: das Ganze besteht, beliebige Grenzen ---
\begin{scope}
\fill[gray!20] (0,0) rectangle (5,2.6);
\draw[thick] (0,0) rectangle (5,2.6);
\draw[dashed] (1.3,0) -- (1.3,2.6);
\draw[dashed] (3.1,0) -- (3.1,2.6);
\draw[dashed] (0,1.1) -- (5,1.1);
\draw[dashed] (0.6,2.6) -- (4.3,0);
\node at (2.5,-0.45) {\small Raum: das Ganze besteht};
\node at (2.5,-0.9) {\small beliebige Grenzen, keine ausgezeichnete};
\end{scope}
% --- Zeit: nur die Grenze besteht ---
\begin{scope}[xshift=7cm]
\draw[dotted,thick,gray] (0,1.3) -- (2.4,1.3);
\draw[dotted,thick,gray] (2.6,1.3) -- (5,1.3);
\fill (2.5,1.3) circle (0.09);
\draw[thick] (2.5,0.9) -- (2.5,1.7);
\node[above] at (2.5,1.75) {\small Jetzt (in recto)};
\node[below] at (1.2,1.15) {\scriptsize vergangen};
\node[below] at (3.8,1.15) {\scriptsize künftig};
\node at (2.5,0.35) {\scriptsize Strecke in obliquo, \zitat{in gar keiner bestimmten Länge}};
\node at (2.5,-0.45) {\small Zeit: nur die Grenze besteht};
\node at (2.5,-0.9) {\small eine ausgezeichnete Grenze, keine beliebige};
\end{scope}
\end{tikzpicture}
\caption{Ganzes und Grenze. Links das topisch Kontinuierliche, das dem Ganzen nach besteht und überall Grenzen zulässt; rechts das chronisch Kontinuierliche, das nur als eine Grenze besteht, deren Strecke ohne Maß mitgemeint ist (nach Brentano). Die beiden Weisen, diskret zu sein, kreuzen sich: ein ausgezeichneter Schnitt ohne beliebige, beliebige Schnitte ohne ausgezeichneten.}
\label{abb:grenze}
\end{figure}

Damit ist die erste Asymmetrie bestimmt. Sie liegt nicht darin, dass die Zeit \zitat{diskret} und der Raum \zitat{kontinuierlich} wäre oder umgekehrt. Beide sind kontinuierlich, und beide lassen sich unterbrechen, aber auf zwei einander ausschließende Weisen. Die Zeit hat ihre Grenze an sich selbst, an genau einer Stelle, und diese Stelle wandert nicht, sondern ist immer dieselbe, das Jetzt, das sich erneuert. Der Raum hat keine Grenze an sich selbst; sie wird ihm gesetzt, an beliebiger Stelle, und keine der Stellen ist vor der anderen ausgezeichnet.

\section{Wo der Raum seine Grenze herhat}

Wenn der Raum von sich aus keine ausgezeichnete Grenze hat, woher hat er dann die, die er offenbar hat? Denn es gibt Orte, und ein Ort ist etwas anderes als eine beliebige Stelle. Die Antwort, die sich aus Hegels Übergang ergibt: Der Ort ist ein als Jetzt gesetzter Punkt. Das ausgezeichnete Hier ist die Zeit, die sich an einer Stelle des Raumes niedergelassen hat. Kochs Antwort ergänzt sie: Das ausgezeichnete Hier ist der Ort des Leibes. Beide Antworten sagen dasselbe, wenn man bedenkt, dass der Leib das Zeitliche ist, das im Raum ist. Er ist nach Brentano die dreidimensionale Grenze eines zeitlichen Kontinuums, das sich in ihm erneuert; er ist nach Koch das Partikulare, das sich a priori selbst individuiert und lokalisiert.

Was also im Raum Grenze ist, im ausgezeichneten Sinn, ist geliehene Zeit. Die Geometrie kennt diese Grenze nicht, und sie braucht sie nicht, denn sie setzt ihre Grenzen überall. Dass es im Raum ein Hier gibt, das kein Dort ist, verdankt sich nicht dem Raum. Das ist das erste Ergebnis, und ihm ist nun in seine beiden Richtungen nachzugehen: zum Hier, das der Leib besetzt, und zum Jetzt, das nach Fink \zitat{weltweit} sein soll.

\chapter{Hier und Jetzt}

\section{Finks Einwand}

Fink hat Hegels sinnlicher Gewissheit einen Fehler vorgeworfen, der genau die Parallelisierung betrifft: Hegel behandle das Hier wie das Jetzt. Beide seien Allgemeine, die sich in jedem Einzelnen verlören. Aber, sagt Fink, das Jetzt ist weltweit. Es ist allen Dingen gemeinsam, es ist die Gleichzeitigkeit, in der alles, was ist, beisammen ist. Das Hier dagegen ist je meines. Es ist der Ort, von dem aus ich sehe, und ein anderes Hier ist für mich ein Dort. Die \zitat{fälschliche Parallelisierung von Hier und Jetzt} verdeckt, dass die Zeit in ihrer Dimension der Gleichzeitigkeit \zitat{den Sinn räumlicher Koexistenz} gewinnt, während der Raum in seinem Hier gar nichts Gemeinsames hat.\footnote{Fink, Frühgeschichte, Kapitel zu Parmenides und zur Weltganzheit.}

Der Einwand ist stark, aber er hat eine merkwürdige Folge. Er macht die Zeit gerade dort zum Gemeinsamen, wo sie am ehesten dem Raum gleicht: im Nebeneinander des Gleichzeitigen. Und er macht den Raum gerade dort zum Eigenen, wo er am wenigsten Raum ist: im Hier, das kein Dort ist. Fink hat also nicht die Zeit gegen den Raum ausgespielt, sondern den raumhaften Aspekt der Zeit gegen den zeithaften Aspekt des Raumes.

\section{Kochs Argument}

Koch nähert sich derselben Frage von der Logik der Einzeldinge her. Es gibt überabzählbar viele reelle Zahlen, aber nur abzählbar viele Sätze; also lassen sich nicht alle Zahlen spezifizieren, und das ist verschmerzbar. Bei Partikularien, den raumzeitlichen Einzeldingen, ist es nicht verschmerzbar: Sie müssen wenigstens teilweise spezifizierbar sein, sonst sind sie keine Gegenstände. Raum und Zeit aber, die Bedingung ihrer Spezifizierbarkeit, sind zugleich ihre Gefahr, denn Raum und Zeit lassen Symmetrien und Wiederholungen zu. Ein spiegelsymmetrisches Universum enthält zu jedem Ding einen Doppelgänger, der ihm qualitativ gleicht und sich von ihm nur numerisch unterscheidet. \zitat{Rechts ist da, wo der Daumen links ist}, sagt man im Scherz; hier wäre es ernst, denn keine Beschreibung könnte das eine vom anderen unterscheiden.\footnote{Koch, Wozu noch Erste Philosophie?, Abschnitt zur Spezifizierbarkeit von Partikularien.}

Kochs Auflösung: Von außen, für eine standpunktneutrale Betrachterin, sind die Doppelgänger unspezifizierbar. Von innen, für einen Betrachter in der Welt, sind sie spezifizierbar durch Indexikalität: ich, hier, jetzt, dieses. Orte und Augenblicke müssen also irreduzibel indexikalischen Charakter haben, und das heißt: Es muss ein Partikulare geben, das sich a priori selbst individuiert und lokalisiert. \zitat{Also sind wir selbst die notwendige Garantie der indexikalischen Spezifizierbarkeit von Partikularien.} Die raumzeitliche Wirklichkeit ist damit nicht subjektiv, aber sie ist \zitat{an ihr selber Erscheinung}: Sie ist so verfasst, dass sie nur von einem Standpunkt in ihr bestimmbar ist.\footnote{Koch, Erste Philosophie, ebd.}

Man sieht sofort, dass Koch das Hier und das Jetzt gleich behandelt. Beide sind indexikalisch, beide kommen von dem Partikularen, das wir sind. Fink würde einwenden: Das Hier ja, das Jetzt nein; das Jetzt ist weltweit, es braucht keinen Standpunkt. Wer hat recht?

\section{Was die Physik entscheidet}

Die Frage lässt sich entscheiden, und zwar an der Stelle, die Reichenbach am genauesten untersucht hat: der Gleichzeitigkeit. Zwei Ereignisse an verschiedenen Orten gleichzeitig zu nennen, setzt voraus, dass man Uhren an verschiedenen Orten verglichen hat. Der Vergleich geschieht durch Signale, und das schnellste Signal ist das Licht. Ein Lichtsignal, das zum Zeitpunkt $t_1$ abgeht, an einem entfernten Ort reflektiert wird und zum Zeitpunkt $t_3$ zurückkehrt, war zu einem Zeitpunkt $t_2$ dort. Welchem? Jeder Zeitpunkt $t_2 = t_1 + \varepsilon\,(t_3 - t_1)$ mit $0 < \varepsilon < 1$ ist zulässig; die Wahl $\varepsilon = \tfrac12$ ist nur durch deskriptive Einfachheit ausgezeichnet, durch die Symmetrie und Transitivität, die sie der Definition verleiht, nicht durch eine Tatsache.\footnote{Reichenbach, Raum-Zeit-Lehre, \S\,19 und \S\,27.} Die Gleichzeitigkeit ferner Ereignisse ist eine Zuordnungsdefinition. Was an ihr Tatsache ist, ist etwas Negatives: \zitat{Gleichzeitigkeit heißt Ausgeschaltetheit des Wirkungszusammenhangs}.\footnote{Reichenbach, Raum-Zeit-Lehre, \S\,22.} Gleichzeitig sind Ereignisse, zwischen denen keine Wirkung laufen kann, und das ist eine ganze Klasse, nicht ein Schnitt.

\begin{figure}[htbp]
\centering
\begin{tikzpicture}[scale=1.05,>=Stealth]
% Achsen
\draw[->] (-3.2,0) -- (3.4,0) node[right] {\small $x$};
\draw[->] (0,-2.6) -- (0,3.2) node[above] {\small $t$};
% Lichtkegel
\fill[gray!12] (0,0) -- (2.6,2.6) -- (-2.6,2.6) -- cycle;
\fill[gray!12] (0,0) -- (2.4,-2.4) -- (-2.4,-2.4) -- cycle;
\draw[gray] (-2.6,-2.6) -- (2.6,2.6);
\draw[gray] (2.6,-2.6) -- (-2.6,2.6);
% Faserbündel: mehrere zeitartige Weltlinien
\foreach \x in {-2.4,-1.6,1.6,2.4} {\draw[thin,gray!70] (\x,-2.4) .. controls (\x+0.15,-0.5) and (\x-0.15,0.8) .. (\x,2.9);}
% eigene Weltlinie (Faser)
\draw[very thick] (0,-2.4) .. controls (0.35,-0.8) and (-0.25,1.2) .. (0.1,2.9);
\fill (0,0) circle (0.07) node[below right] {\small Jetzt};
% Gleichzeitigkeitsflächen mit verschiedenen epsilon
\draw[dashed] (-3,-0.7) -- (3,0.7);
\draw[dashed] (-3,0) -- (3,0);
\draw[dashed] (-3,0.7) -- (3,-0.7);
\node[right] at (3.05,0.75) {\scriptsize $\varepsilon\neq\tfrac12$};
\node[right] at (3.05,-0.75) {\scriptsize $\varepsilon\neq\tfrac12$};
\node[right] at (3.05,0.15) {\scriptsize $\varepsilon=\tfrac12$};
% Beschriftungen
\node at (0,2.3) {\small früher $\to$ später};
\node[align=center] at (-1.9,1.5) {\scriptsize zeitfolge-\\[-2pt]\scriptsize unbestimmt};
\node[align=center] at (1.9,-1.5) {\scriptsize zeitfolge-\\[-2pt]\scriptsize unbestimmt};
\node[below] at (0.9,-2.55) {\small Nebeneinander der Fasern};
\node[right] at (0.2,2.75) {\small Faserrichtung};
\end{tikzpicture}
\caption{Faser und Nebeneinander. Die dick gezeichnete Weltlinie ist die Faser, längs deren dieselbe Materie ihre Zustände durchläuft (Genidentität); die dünnen Linien sind die Fasern anderer Dinge, deren Nebeneinander der Raum ist. Die gestrichelten Geraden sind mögliche Gleichzeitigkeitsflächen durch das Jetzt; welche von ihnen gilt, ist Definition ($\varepsilon$), was zwischen ihnen liegt, ist die Klasse der Ereignisse ohne Wirkungszusammenhang (nach Reichenbach).}
\label{abb:faser}
\end{figure}

Das entscheidet nicht zwischen Fink und Koch, aber es grenzt ab, worüber zu entscheiden ist. Fink meint mit dem weltweiten Jetzt keine durch Uhrenvergleich definierte Beziehung, sondern das Anwesen, in dem alles, was ist, versammelt ist, und er bestimmt es ausdrücklich nicht vom Jetzt-Sagen der Seele her. Die Physik trifft dieses Jetzt nicht. Sie trifft nur seine Verlängerung ins Bestimmbare: Sobald man sagen will, welche fernen Ereignisse im weltweiten Jetzt beisammen sind, ist man bei der Gleichzeitigkeit, und die ist Definition. Was dem weltweiten Jetzt dann bleibt, hat Fink selbst benannt: In ihrer Dimension der Gleichzeitigkeit gewinnt die Zeit den Sinn räumlicher Koexistenz. Das weltweite Jetzt ist die Zeit in der Gestalt des Nebeneinander, und an dieser Gestalt hat die Physik gezeigt, dass sie keinen Schnitt enthält, sondern eine Klasse. Was an der Zeit keine Klasse ist, ist die Ordnung längs der Weltlinie: Zwischen zwei Ereignissen an demselben Ding gibt es ein Früher und Später, das keine andere Wahl der Gleichzeitigkeit umkehrt. Das Jetzt, das nicht in eine Klasse zerfließt, ist also das Jetzt an der Faser, und die Faser ist die Weltlinie eines Dinges. Das ist nicht Kochs Argument, aber es trifft sich mit seinem Ergebnis: Das Jetzt, das etwas auszeichnet, gehört zu einem Partikularen. Koch selbst nimmt für sein symmetrisches Universum keinen objektiven Zeitpfeil an; die Physik gibt dem indexikalischen Jetzt etwas hinzu, das bei ihm nicht vorkommt, die Eigenzeit des Dinges, an dem es hängt.

\section{Das Hier, noch einmal}

Umgekehrt lässt sich nun auch das Hier genauer fassen. Finks Wort, das Hier sei mein Ort, ist richtig, aber es sagt nicht, woher das Hier seine Auszeichnung hat. Der Raum, als Nebeneinander der Fasern, kennt keinen ausgezeichneten Punkt; das war das Ergebnis an der Grenze. Das Hier ist der Punkt, an dem eine Faser den Raum durchstößt, und zwar diejenige Faser, längs deren ich mich erneuere. Das Hier ist also die Spur einer Zeit im Raum. Die Lesart des Ortes als gesetztes Jetzt und Kochs Formel vom Leib als der Garantie der Spezifizierbarkeit sagen beide dies.

Das Ergebnis ist eine Kreuzung, wie schon bei der Grenze. Das Jetzt, das allen gemeinsam ist, ist Raum. Das Hier, das nur meines ist, ist Zeit. Die Parallelisierung von Hier und Jetzt, die Fink Hegel vorwirft, ist deshalb nicht bloß falsch; sie ist die Verwechslung zweier Kreuzungen, die einander spiegeln. Die Asymmetrie tritt genau dann hervor, wenn man in jedem der beiden fragt, was an ihm vom anderen stammt.

\chapter{Richtung}

\section{Die Faserrichtung}

Reichenbach baut die Ordnung der Zeit nicht auf ein Zeitgefühl, sondern auf die Kausalität. Sein Kennzeichenprinzip lautet: Wenn ein Ereignis $E_1$ Ursache eines Ereignisses $E_2$ ist, dann erscheint eine kleine Veränderung an $E_1$ auch an $E_2$, aber nicht umgekehrt. Man kann das Kennzeichen an der Ursache anbringen und es an der Wirkung wiederfinden; an der Wirkung angebracht, findet man es an der Ursache nicht. Diese Asymmetrie braucht keinen Zeitbegriff, sie definiert ihn. Früher ist, was das Kennzeichen abgibt, später, was es empfängt.\footnote{Reichenbach, Raum-Zeit-Lehre, \S\,21.}

Auf dieser Grundlage wird die Zeit zu dem, was der Raum nicht ist: \zitat{Eine hinzukommende wesentliche Eigenschaft der Zeit ist die Auszeichnung eines Richtungssinnes; und dieser wieder beruht darauf, daß die Zeit --- und nur sie --- die Dimension der Wirkungslinien ist [\ldots]. So ist die Zeit die Faserrichtung der Mannigfaltigkeit, diejenige Richtung, in der sich die Kausalketten erstrecken, während der Raum nur das Nebeneinander der Faserbündel aufnimmt.}\footnote{Reichenbach, Raum-Zeit-Lehre, \S\,43.} Ein ausgezeichneter Spezialfall der Kausalkette ist die Weltlinie eines Dinges, das mit sich identisch bleibt; Reichenbach nennt diese Beziehung Genidentität. Längs der Zeit liegen Zustände desselben Dinges, längs des Raumes liegen verschiedene Dinge. Das ist die zweite Asymmetrie in ihrer physikalischen Form: Die Zeit verbindet, was eines ist, der Raum trennt, was vieles ist.

Was Reichenbach nicht löst, ist die Einsinnigkeit. Dass die Zeit \emph{fortschreitet}, dass es einen Unterschied macht, in welcher der beiden Richtungen die Faser durchlaufen wird, ist durch das Kennzeichenprinzip zwar definierbar, aber die Definition trifft nicht das Gemeinte. \zitat{Fortschreiten der Zahlen heißt nichts anderes als Richtung vom kleineren zum größeren. Mit Fortschreiten der Zeit meinen wir aber etwas Selbständiges, ein evidentes Erlebnis und zugleich eine Aussage über die Wirklichkeit.}\footnote{Reichenbach, Raum-Zeit-Lehre, \S\,21.} Das Problem wird vertagt und im Buch nicht wieder aufgenommen. Es bleibt auch hier offen; es wird aber zuletzt deutlich werden, warum es nicht auf der Ebene der Physik lösbar ist.

\section{Der Raum hat keine Richtung}

Dass der Raum aus sich keine Richtung hat, ist eine alte Einsicht mit vielen Zeugen. Hegel sagt, die drei Dimensionen des Raumes, Höhe, Länge und Breite, seien nur der Zahl nach unterschieden; welche man Höhe nennt, ist gleichgültig. Bonsiepen hat gezeigt, dass Hegel dabei Aristoteles vor Augen hat, der erklärt, nicht in jedem Körper sei ein Oben und Unten, Rechts und Links, Vorn und Hinten zu finden, sondern nur in beseelten Körpern; die Richtungen sind mit dem Ursprung der Bewegung verknüpft.\footnote{Bonsiepen, Hegels Raum-Zeit-Lehre, zu Enzyklopädie \S\,255 und De caelo II.} Martić hält an Newton dieselbe Sache fest: Newton unterscheidet die Anisotropie der Zeit, die eine feste Folge aufweist, von der Isotropie des Raumes, dessen Teile keine Reihenfolge haben.\footnote{Martić, Zeit und Raum, Kapitel zu Newton.} Und Kant hat 1768 an den inkongruenten Gegenstücken, der linken und der rechten Hand, gezeigt, dass die Unterscheidung von rechts und links durch keine Beschreibung der inneren Verhältnisse eines Körpers zu gewinnen ist; er schloss daraus auf einen absoluten Raum, den er später, in den Metaphysischen Anfangsgründen, als bloße Idee von einem Raume bestimmte, der in der Tat kein Raum ist.\footnote{Schneider, Raum-Zeit-Problem, Kapitel zum absoluten Raum; Martić, Zeit und Raum, zu den Gegenden im Raume.}

Kochs Argument vom symmetrischen Universum ist die Zuspitzung: In einem symmetrischen Universum ist das rechte Ding vom linken nicht zu unterscheiden, außer für einen, der drinsteht. Der Raum, das reine Nebeneinander, kennt weder ein Rechts noch ein Oben noch ein Vorn.

\begin{figure}[htbp]
\centering
\begin{tikzpicture}[scale=1,>=Stealth]
% Spiegelachse
\draw[thick,dashed] (0,-2.2) -- (0,2.6);
\node[above] at (0,2.6) {\small Spiegelebene};
% Linke Figur
\begin{scope}[xshift=-2.4cm]
\draw[thick] (0,0.8) circle (0.35);
\draw[thick] (0,0.45) -- (0,-0.6);
\draw[thick] (0,0.2) -- (-0.7,-0.2);
\draw[thick] (0,0.2) -- (0.7,0.4);
\draw[thick] (0,-0.6) -- (-0.4,-1.4);
\draw[thick] (0,-0.6) -- (0.4,-1.4);
\node[below] at (0,-1.5) {\small $A$};
\draw[->,thick] (0.75,0.4) -- (1.3,0.4) node[right] {\scriptsize \zitat{rechts}};
\end{scope}
% Rechte Figur (gespiegelt)
\begin{scope}[xshift=2.4cm]
\draw[thick] (0,0.8) circle (0.35);
\draw[thick] (0,0.45) -- (0,-0.6);
\draw[thick] (0,0.2) -- (0.7,-0.2);
\draw[thick] (0,0.2) -- (-0.7,0.4);
\draw[thick] (0,-0.6) -- (0.4,-1.4);
\draw[thick] (0,-0.6) -- (-0.4,-1.4);
\node[below] at (0,-1.5) {\small $A'$};
\draw[->,thick] (-0.75,0.4) -- (-1.3,0.4) node[left] {\scriptsize \zitat{rechts}};
\end{scope}
% Außenbetrachter
\node[align=center] at (0,-2.9) {\small von außen: $A$ und $A'$ qualitativ gleich, nur numerisch verschieden\\[-1pt]\small von innen: \zitat{hier} und \zitat{dort}, \zitat{ich} und \zitat{er}};
\end{tikzpicture}
\caption{Das symmetrische Universum. Zu jedem Ding ein Doppelgänger, dem keine Beschreibung ein Merkmal geben kann, das der andere nicht auch hätte. Die Unterscheidung gelingt nur indexikalisch, von einem Standpunkt in der Welt (nach Koch). Der Raum selbst liefert kein Rechts und kein Links.}
\label{abb:doppelgaenger}
\end{figure}

\section{Woher die Richtungen des Raumes kommen}

Wenn der Raum keine Richtungen hat, aber die Erfahrung sie kennt, dann müssen sie von woanders kommen. Koch gibt in seiner Lektüre der Hegelschen Naturphilosophie eine Herleitung, die genauer ist, als man erwartet. Die Richtungen des Raumes kommen vom Leib, und zwar von drei Weisen, wie der Leib sich bewegt. Die natürliche Bewegung, der Fall, gibt Oben und Unten. Die freie Bewegung, das Gehen, gibt Hinten und Vorn: Vorn ist, wohin ich gehe. Und in der freien Bewegung gabelt sich der Weg: Die Wahl zwischen den beiden Zweigen gibt Rechts und Links. Die Modi der Zeit dagegen, Vergangenheit, Gegenwart und Zukunft, kommen von den Seelenvermögen: Erinnerung, Wahrnehmung, Erwartung.\footnote{Koch, Hegelsche Subjekte in Raum und Zeit.} Der Leib ist also nicht nur ein Ding im Raum, das zufällig Richtungen hat; er ist das, was dem Raum seine Richtungen \emph{gibt}, indem er sich in ihm bewegt.

Baumann hat behauptet, dass Bestimmungen wie rechts und links, oben und unten \zitat{Eigenthum des Geistes} sind und im Körper nur mit angedeutet werden.\footnote{Baumann, Lehren, Bd.\,2, S.\,657, Kurzer Lehrbegriff.} Das ist dieselbe Einsicht in älterer Sprache: Die Richtung ist nicht im Raum, sondern kommt zu ihm hinzu, von dem, was in ihm lebt.

Nun ist die Bewegung, aus der Koch die Richtungen herleitet, ein Übergang, und der Übergang war schon bei Zenon das, was der Raum zerlegt und die Zeit trägt. Man darf daraus nicht schließen, die Richtungen des Raumes kämen aus der Zeit; Fink hat mit Recht darauf bestanden, dass Bewegung Räumen und Zeitigen zugleich ist, und Fallen und Gehen sind so räumlich wie zeitlich. Was sich schließen lässt, ist weniger und fester: Die Richtungen des Raumes kommen nicht aus dem Raum. Sie kommen von einem Bewegten, und ein Bewegtes ist, was längs einer Faser dasselbe bleibt und dabei seinen Ort wechselt. Die zweite Asymmetrie lautet also: Die Zeit hat eine Richtung an sich selbst, weil sie die Faser ist; der Raum hat Richtungen nur, weil eine Faser ihn durchläuft.

\section{Auch die Zahl der Dimensionen}

Es gibt eine Stelle, an der sich diese Herkunft überraschend weit treiben lässt: die Dimensionszahl. Reichenbach fragt, was es heißt, dass der Raum drei Dimensionen hat, wenn doch jede Mannigfaltigkeit von Ereignissen sich durch beliebig viele Parameter beschreiben lässt. Sein Beispiel: Die Bewegung zweier Punkte auf einer Geraden lässt sich ebenso gut als Bewegung eines Punktes in einer Ebene auffassen. Beide Beschreibungen sind äquivalent, solange man nur die Lagen betrachtet. \zitat{Nehmen wir aber das Nahwirkungsprinzip hinzu, so verschwindet die Äquivalenz.} Eine Störung, die sich von einem Punkt aus stetig ausbreitet, breitet sich in der einen Darstellung stetig aus und in der anderen längs zweier gekreuzter Geraden mit unendlicher Geschwindigkeit. \zitat{Das Nahwirkungsprinzip kann nur für eine einzige Wahl der Dimensionszahl des Parameterraums erfüllt sein; derjenige Parameterraum, in dem es gilt, heißt Koordinatenraum oder ›wirklicher Raum‹.}\footnote{Reichenbach, Raum-Zeit-Lehre, \S\,44.}

\begin{figure}[htbp]
\centering
\begin{tikzpicture}[scale=1,>=Stealth]
% Links: 1-dim Koordinatenraum mit zwei Punkten und Störung
\begin{scope}
\draw[thick] (-0.3,0) -- (4.3,0);
\fill (1,0) circle (0.07) node[below] {\scriptsize $p_1$};
\fill (3,0) circle (0.07) node[below] {\scriptsize $p_2$};
\draw[->,thick,gray] (1,0.25) -- (1.7,0.25);
\draw[->,thick,gray] (1,0.25) -- (0.3,0.25);
\node[above] at (1,0.35) {\scriptsize Störung von $p_1$};
\node at (2,-0.8) {\small Nahwirkung stetig};
\node at (2,-1.2) {\small \zitat{wirklicher Raum}};
\end{scope}
% Rechts: 2-dim Parameterraum, Punkt (x1,x2), Störung greift auf gekreuzten Geraden
\begin{scope}[xshift=6.5cm]
\draw[->] (-0.3,0) -- (3.6,0) node[right] {\scriptsize $x_1$};
\draw[->] (0,-0.3) -- (0,3.4) node[above] {\scriptsize $x_2$};
\fill (1,3) circle (0.07) node[above right] {\scriptsize $P=(x_1,x_2)$};
\draw[gray,thick] (0.2,3) -- (3.3,3);
\draw[gray,thick] (1,0.2) -- (1,3.3);
\draw[->,thick,gray] (1,3) -- (1.6,3);
\draw[->,thick,gray] (1,3) -- (0.4,3);
\draw[->,thick,gray] (1,3) -- (1,2.4);
\node[align=center] at (1.7,-0.8) {\small dieselbe Störung: sogleich auf\\[-1pt]\small zwei gekreuzten Geraden};
\end{scope}
\end{tikzpicture}
\caption{Parameterraum und Koordinatenraum. Zwei Punkte auf einer Geraden oder ein Punkt in der Ebene: Als Lagen sind beide Darstellungen gleichwertig. Erst die Forderung, dass Wirkungen sich stetig ausbreiten, zeichnet eine aus. Die Dimensionszahl des Raumes ist der objektive Sinn des Nahwirkungsprinzips (nach Reichenbach).}
\label{abb:parameterraum}
\end{figure}

Das heißt: Sogar die Zahl der Dimensionen des Raumes ist der Sinn eines Wirkungsgesetzes. Wirkung aber läuft längs der Faser, in der Zeit. Kant hat 1747, in den Gedanken von der wahren Schätzung der lebendigen Kräfte, dieselbe Vermutung geäußert: Die Dreidimensionalität des Raumes folge aus dem Gesetz, nach dem die Substanzen aufeinander wirken, dem umgekehrten Quadrat der Entfernung; hätten die Substanzen keine Kraft, außer sich zu wirken, so gäbe es keinen Raum.\footnote{Schneider, Raum-Zeit-Problem, Kapitel zu Kants vorkritischen Schriften.} Der junge Kant und der späte Reichenbach sagen dasselbe: Was am Raum Struktur ist, bis hinab zur Zahl seiner Dimensionen, wird von der Wirkung bestimmt, die in ihm läuft, und die Wirkung ist zeitlich.

\chapter{Maß und Unmaß}

\section{Weile und Weite}

Ute Guzzoni hat dem Wort \zitat{Weile} nachgefragt, das eine Phase der Zeit nennt, ohne sie zu messen. Eine Weile ist nicht eine Abfolge von Jetztpunkten, sondern Gegenwärtigkeit ohne Verlauf, erfüllte Zeit. Ihre Dauer ist nur qualitativ bestimmt; es gibt keine halbe Weile und keine doppelte. Grimm hat es gewusst: die Formel \zitat{eine weile} macht ihr Glück durch ihre Unbestimmtheit. Und die Weile hat, gegen den Fluss der Zeit gehalten, etwas von einer Fläche: \zitat{Diesem Fluß der Zeit gegenüber erscheint die Weile eher als eine Fläche oder ein Raum, sie verläuft gerade nicht linear.}\footnote{Guzzoni, Weile und Weite, Kapitel zur Weile.}

Das Gegenstück ist die Weite. Höhe, Breite und Länge sind messbar; die Weite ist es nicht. Die Weite ist der Raum, wie er erfahren wird, wenn er nicht vermessen wird: das Offene, in das man hinausschaut, die Gegend, in der man sich aufhält. Guzzoni stellt beide zusammen und findet in ihnen ein ursprüngliches Ineinander von Zeit und Raum, das vor jeder Metrik liegt. Der Ort ist der Augenblick des Raumes; die Gegend entspricht der Weile.\footnote{Guzzoni, Weile und Weite, Kapitel zu Ort und Gegend.} Das sind keine Bilder. Es ist die Erfahrung, dass die Zeit, wenn sie weilt, Raum gewinnt, und dass der Raum, wenn er Gegend ist, Zeit hat.

Die Richtung dieses Übergangs ist bemerkenswert. Heidegger: \zitat{Die Zeit wird zum Raum.} Wagner: \zitat{Zum Raum wird hier die Zeit.} Hölderlin: \zitat{Zeit eilt hin zum Ort.} Guzzoni versammelt diese Sätze und bemerkt, dass sie alle in dieselbe Richtung gehen, und nicht umgekehrt; niemand sagt, der Raum werde zur Zeit. \zitat{In der Erfahrung scheint die Zeit das Vorrangige und Bestimmende zu sein. Und doch ist sie immer schon eingeräumt, mit ihrem Ablauf wie mit ihrem Weilen und Verweilen skandiert sie den Raum, in den sie hineingehört.}\footnote{Guzzoni, Weile und Weite, Kapitel zu Heidegger.} Das Unmaß, die Weile und die Weite, ist beiden gemeinsam und geht dem Maß voraus. Aber im Unmaß hat die Zeit den Vorrang und der Raum die Aufnahme. Der Raum räumt ein, was die Zeit bringt.

\section{Topologie vor Metrik}

Reichenbach kommt von der entgegengesetzten Seite, von der Messung, und findet dieselbe Reihenfolge. Die Metrik des Raumes beruht auf einer Zuordnungsdefinition: Man muss festlegen, welcher Körper als starr gilt, welche Kraft als universell und deshalb wegzudefinieren. Die Geometrie, die sich ergibt, ist die Geometrie plus die Festsetzung; beides zusammen ist empirisch, keines allein.\footnote{Reichenbach, Raum-Zeit-Lehre, \S\,3--\S\,8.} Vor der Metrik liegt die Topologie: die Nachbarschaftsverhältnisse, das Zwischen, die Ordnung. Und die Topologie des Raumes wird, wie sich an der Dimensionszahl zeigte, durch das Nahwirkungsprinzip bestimmt, also durch die Kausalität, also durch die Zeit. Die Metrik kommt zuletzt, und sie kommt bei Reichenbach aus der Gravitation: Das metrische Feld ist das Gravitationsfeld, die Geometrie wird zum Ausdruck der Schwere.\footnote{Reichenbach, Raum-Zeit-Lehre, \S\,38--\S\,40.}

Für die Zeit gilt dasselbe. Die Metrik der Zeit, welche Uhr gleichförmig läuft, ist Zuordnungsdefinition, und auch die Ordnung der Zeit beruht auf einer, die Reichenbach die topologische Zuordnungsdefinition der Zeitfolge nennt: Ist $E_2$ die Wirkung von $E_1$, so heißt $E_2$ später als $E_1$.\footnote{Reichenbach, Raum-Zeit-Lehre, \S\,21.} Die Asymmetrie liegt also nicht darin, dass die Zeit ohne Definition auskäme. Sie liegt in der Reihenfolge der Definitionen. Die Zeitfolge wird an der Kausalkette definiert und sonst an nichts; die Ordnung des Raumes, das Zwischen und die Nachbarschaft, wird an den Kausalketten definiert, also über die Zeit. Und die Definition der Zeitfolge lässt keine Wahl offen, sobald das Kennzeichen läuft, während die Gleichzeitigkeit eine Wahl offen lässt, den Wert von $\varepsilon$. Beide haben eine definitorische Metrik, beide eine definitorische Ordnung; aber die eine Ordnung wird an der anderen definiert, nicht umgekehrt.

\section{Das Lichtjahr}

Reichenbach spricht die Ableitung aus, zuerst bei der Lichtgeometrie: Die Raummessung \zitat{läßt sich zurückführen auf die Zeitmessung; die Zeit ist gegenüber dem Raum das logisch Primäre}.\footnote{Reichenbach, Raum-Zeit-Lehre, \S\,27.} Und dann, bei der Raum-Zeit-Ordnung, in aller Schärfe: Das \zitat{›Lichtjahr‹ der Astronomen, ursprünglich aus reinen Zweckmäßigkeitsgründen eingeführt, trifft gerade die logische Urform aller Längenmessung. Die Zeit und durch sie die Kausalität liefert das Maß und die Ordnung des Raumes}.\footnote{Reichenbach, Raum-Zeit-Lehre, \S\,42.}

Wie das geht, zeigt ein einziges Verfahren. Von einer Uhr geht ein Lichtsignal aus, wird an einem entfernten Körper reflektiert und kehrt zurück. Die Uhr misst die Zeit $\Delta t$ zwischen Abgang und Rückkehr. Die Entfernung des Körpers ist dann $\tfrac12\, c\, \Delta t$, wobei $c$ nicht gemessen, sondern gesetzt wird, denn die Lichtgeschwindigkeit ist die Einheit, in der Längen in Zeiten übersetzt werden. Man braucht zur Längenmessung also nur eine Uhr und das Licht. Man braucht keinen Maßstab.

\begin{figure}[htbp]
\centering
\begin{tikzpicture}[scale=1,>=Stealth]
% Weltlinie der Uhr
\draw[very thick] (0,-0.3) -- (0,4.6);
\node[above] at (0,4.6) {\small Uhr};
% Weltlinie des fernen Körpers
\draw[thick,gray] (3,-0.3) -- (3,4.6);
\node[above] at (3,4.6) {\small ferner Körper};
% Signal hin und zurück
\draw[->,thick] (0,0.5) -- (3,2.2);
\draw[->,thick] (3,2.2) -- (0,3.9);
\fill (0,0.5) circle (0.07) node[left] {\small $t_1$};
\fill (3,2.2) circle (0.07) node[right] {\small $t_2$};
\fill (0,3.9) circle (0.07) node[left] {\small $t_3$};
% Zeitintervall
\draw[<->] (-1.1,0.5) -- (-1.1,3.9);
\node[left] at (-1.1,2.2) {\small $\Delta t = t_3 - t_1$};
% Länge
\draw[<->,dashed] (0,-0.15) -- (3,-0.15);
\node[below] at (1.5,-0.2) {\small $\ell = \tfrac12\,c\,\Delta t$};
\node[align=center] at (1.5,-1.1) {\small Länge aus Zeit: kein Maßstab, nur Uhr und Licht};
\end{tikzpicture}
\caption{Das Lichtjahr als Urform der Länge. Ein Signal geht bei $t_1$ ab, wird reflektiert und kehrt bei $t_3$ zurück. Die Entfernung ist die halbe Laufzeit, in Lichteinheiten. Was der Raum an Maß hat, wird von einer Uhr abgelesen (nach Reichenbach).}
\label{abb:lichtjahr}
\end{figure}

Reichenbach selbst setzt dem Verfahren eine Grenze, und sie muss mitgenannt werden. Die Lichtgeometrie erreicht nur eine Klasse von Raumsystemen, die sich durch die Wahl des Zeitparameters unterscheiden; um unter ihnen das eine auszuzeichnen, das der Raum der Körper ist, braucht man einen materiellen Messkörper. Der Raum gibt sich also nicht restlos aus der Zeit her. Es bleibt ein Rest, und der Rest ist der Körper, dessen Starrheit man voraussetzt. Was das für den Raum bedeutet, wird zuletzt zu bedenken sein.

\section{Der Abstand ohne Unterschied}

Brentano hat eine dritte Weise entdeckt, in der das Maß der Zeit anders ist als das des Raumes. Zwei Punkte im Raum, die voneinander entfernt sind, unterscheiden sich sachlich: Hier ist nicht Dort, so wie Rot nicht Blau ist. Der räumliche Abstand ist eine Differenz des Gegenstandes. Zwei Punkte in der Zeit, die voneinander entfernt sind, brauchen sich sachlich nicht zu unterscheiden. Eine Melodie kann wiederkehren, ein Ton kann derselbe sein, ein Schlaf kann ohne jede Veränderung vergehen. Der zeitliche Abstand ist keine Differenz des Gegenstandes, sondern eine Differenz der Weise, wie derselbe Gegenstand anerkannt wird. Deshalb ist die Raumanschauung ein Kontinuum von Objekten, die Zeitanschauung ein Kontinuum von Modi.\footnote{Brentano, Untersuchungen, Kapitel VIII, Das Zeitliche als Relatives.}

Das Maß der Zeit misst also nichts am Gemessenen. Es misst den Abstand zwischen zwei Anerkennungen desselben. Das ist der Grund, warum es keine Wissenschaft der Zeit gibt, die der Geometrie entspräche. Hegel hat es gesagt: Die Arithmetik, die man der Zeit zuordnen möchte, reduziert sie auf das tote Eins; die Zeit lässt sich nicht so konstruieren, wie die Geometrie den Raum konstruiert.\footnote{Bonsiepen, Hegels Raum-Zeit-Lehre, zu Enzyklopädie \S\,259.} Und Aristoteles hat die Zeit als Zahl der Bewegung bestimmt, nicht als Größe: Die Zahl umfasst die Bewegung, wie das Gefäß den Körper umfasst, aber sie ist nicht dessen Ausdehnung. Fink hält fest, dass bei Aristoteles der Raum endlich, die Zeit un-endlich ist, und dass erst das Leere, das Kenon, zwischen beiden vermittelt; von ihm her erst werde die Fundierung der Zeit in der Größe ganz einsichtig.\footnote{Fink, Frühgeschichte, Kapitel zur aristotelischen Zeit-Analytik.}

\section{Ergebnis}

Die dritte Asymmetrie hat zwei Seiten. Auf der Seite des Maßes kommt das Maß des Raumes aus der Zeit: Die Ordnung des Raumes wird an den Kausalketten definiert, seine Metrik wird von Uhren abgelesen, und was er darüber hinaus an Maß hat, hat er von einem Körper, den man als starr voraussetzt. Auf der Seite des Unmaßes, der Weile und der Weite, sind Zeit und Raum gemeinsam vor aller Metrik; aber auch dort ist die Zeit das Bewegende und der Raum das Aufnehmende, die Zeit wird zum Raum, nicht der Raum zur Zeit.

Hier ist ein Einwand zu bedenken, der mit Kant wiederkehren wird. Auch die Uhr ist ein Körper, und ohne ein Beharrliches wird die Zeit nicht zur Größe. Misst die Uhr den Raum dann wirklich aus der Zeit, oder aus einem räumlichen Beharrlichen? Die Antwort liegt in dem, was die Uhr braucht. Sie braucht ein Ding, das längs seiner Weltlinie dasselbe bleibt und sich periodisch erneuert; sie braucht keinen Maßstab, keine Länge, keine Raummetrik. Die Lichtgeometrie braucht die Uhr und das Signal, und für den letzten Schritt den starren Körper. Das Zeitmaß setzt also ein Beharrliches voraus, aber kein Raummaß; das Raummaß setzt das Zeitmaß voraus und außerdem ein Beharrliches. Das Beharrliche ist im Raum, aber es ist nicht der Raum. Das ist die Asymmetrie des Maßes, und der Rest, den der Raum nicht hergibt, hat damit einen Namen: das Ding, das beharrt.

\chapter{Sachlich und modal}

\section{Differenzen des Gegenstandes, Differenzen des Anerkennens}

Die drei Asymmetrien, Grenze, Richtung und Maß, haben einen gemeinsamen Grund, und Brentano hat ihn benannt. Räumliche Differenzen sind Differenzen des Gegenstandes. Was hier ist, ist ein anderes als was dort ist, so wie Rot ein anderes ist als Blau. Zeitliche Differenzen sind keine Differenzen des Gegenstandes. Dass etwas vergangen, gegenwärtig oder künftig ist, ist keine Bestimmung an ihm, sondern eine Weise, wie es vorgestellt und anerkannt wird: ein Modus. Vergangen und künftig sind Modi in obliquo, gegenwärtig ist der Modus in recto. Und diese Modi sind komparative Relativa zum Sprechenden: Was für mich vergangen ist, ist es im Verhältnis zu meinem Jetzt.\footnote{Brentano, Untersuchungen, Kapitel VII, Sachliche und modale Differenzen, und VIII, Das Zeitliche als Relatives.}

Brentano prüft, was geschähe, wenn es anders wäre. Gäbe es nur einen einzigen Präteritalmodus, so wäre das Vergangene vom Gegenwärtigen unterschieden wie ein Ort vom anderen; die Zeit wäre dann \zitat{ein Topoid von einer Dimension}, ein eindimensionaler Raum, und das Nacheinander wäre ein Nebeneinander.\footnote{Brentano, Untersuchungen, Kapitel IV, zum Präteritalmodus.} Dass die Zeit kein Topoid ist, liegt daran, dass die Modi kontinuierlich variieren, und dass sie eine einzige Dimension hat, ist keine Zufälligkeit wie beim Raum: Ein Chronoid von anderer Dimensionszahl ist \zitat{unmittelbar absurd}, während Topoide von beliebiger Dimensionszahl denkbar sind. Der Raum könnte anders sein, als er ist; die Zeit nicht.

Die Überordnung, die daraus folgt, spricht Brentano aus: Der Begriff des Zeitlichen \zitat{verhält sich zu dem des Räumlichen als ein übergeordneter} und begreift Spezies unter sich, die aller räumlichen Kontinuität ermangeln.\footnote{Brentano, Untersuchungen, S.\,213.} Alles Reale ist zeitlich; nicht alles Reale ist räumlich. Und noch einmal die Kreuzung, jetzt an Hier und Jetzt: \zitat{Der Ortspunkt ist Grenze für ein Räumliches in vielen, ja vielleicht in allen denkbaren räumlichen Richtungen, aber niemals Grenze in einer der beiden zeitlichen Richtungen, und vom Zeitpunkt gilt das Umgekehrte.}\footnote{Brentano, Untersuchungen, S.\,208.}

Brentano hat diese Lehre später abgeschwächt und auch in der Zeit relative sachliche Differenzen zugelassen, den Abstand als Bestimmung des Gegenstandes. Aber der Kern bleibt: Der Abstand in der Zeit ist ein Abstand zwischen Weisen, in denen derselbe Gegenstand anerkannt wird. Man kann sich denselben Ton, dieselbe Farbe, denselben Zustand nach beliebig langer Zeit wiederholt denken, ohne dass sich etwas an ihm geändert hätte. Man kann sich nicht denselben Punkt an zwei Orten denken.

\section{Warten und Gehen}

Chiba gibt der modalen Asymmetrie eine Form, die ohne Brentanos Terminologie auskommt: \zitat{Um die Zukunft zu erfahren, muss man nur warten. Dergleichen gilt nicht für den Raum.}\footnote{Chiba, Kants Ontologie der raumzeitlichen Wirklichkeit, Kapitel zur raumzeitlichen Wirklichkeit.} Um das Dort zu erfahren, muss man hingehen, und Gehen ist eine Tätigkeit, die im Raum eine Richtung wählt. Warten ist keine Tätigkeit, die eine Richtung wählt. Die Zukunft kommt von selbst; das Dort kommt nicht von selbst. Das ist die Einsinnigkeit, die Reichenbach nicht definieren konnte, in der Sprache der Erfahrung. Sie ist kein Merkmal der Ereignisse, sondern der Weise, wie sie sich einstellen.

Chiba fasst die raumzeitlichen Gegenstände als solche, die räumlich \emph{und/oder} zeitlich sind, und die Zeit ist die weitere Klasse: Es gibt Zeitliches ohne Räumliches, ein Gedanke etwa, aber kein Räumliches ohne Zeitliches. Das ist Brentanos Überordnung bei Kant. Und Chiba gibt dem Realismus, den er bestreitet, einen Namen, der die Sache trifft: Er sei ein \zitat{zeit-transzendenter Standpunkt}, ein Blick von nirgendwann. Die Existenz eines Gegenstandes ist für ihn dessen Erkennbarkeit; leere Zeit und leerer Raum sind nichts Wirkliches, weil sie nichts sind, was erkannt werden könnte.\footnote{Chiba, Kants Ontologie, Kapitel zum Anti-Realismus und zum Existenzbegriff.}

Der Ausdruck \zitat{zeit-transzendent} ist genauer, als er zunächst scheint. Der Realismus, der die Welt als das nimmt, was sie unabhängig von jedem Standpunkt ist, muss sich einen Blick denken, der nicht jetzt ist. Einen Blick, der nicht hier ist, kann man sich leicht denken: Man geht einfach woanders hin, oder man denkt sich alle Orte zugleich. Einen Blick, der nicht jetzt ist, kann man sich nur denken, indem man die Zeit zum Raum macht, in dem alle Jetzt nebeneinander liegen. Der zeit-transzendente Standpunkt ist also der Standpunkt, von dem aus die Zeit als Raum erscheint.

\section{Die Physik und das Nichts}

Koch hat dieselbe Sache von der Physik her gesagt. Die Physik abstrahiert von der Subjektivität. Sie hat die Kontinuität der Zeit mittels analytischer Geometrie in diskrete reelle Zahlen aufgelöst; die Zeit wird zum Parameter $t$, der wie eine Koordinate behandelt wird. Damit hat sie genau das von der Zeit behalten, worin die Zeit dem Raum gleicht, die Ordnung, und das weggelassen, worin sie sich unterscheidet, den Modus. Koch nennt die Subjektivität das Nichts, den Ursprung der Negation, des Werdens, der Bewegung, und die Physik den Versuch, die Welt unter Absehung vom Nichtigen zu begreifen, ein eleatisches Unternehmen.\footnote{Koch, Erste Philosophie, Abschnitt zur Physik und zur Subjektivität.}

Das ist keine Kritik an der Physik. Es ist die Feststellung, was sie tut. Reichenbach selbst hat es festgestellt, als er die Wahrnehmungsordnung von der Ereignisordnung unterschied: Die Wahrnehmungen bilden eine eindimensionale Kette, die objektiven Ereignisse eine Mannigfaltigkeit mit einer Dimension mehr, in der nichts \zitat{jetzt} ist.\footnote{Reichenbach, Raum-Zeit-Lehre, \S\,45.} Und er hat die Ich-Zeit, in der das \zitat{ewige Jetzt} vorkommt, ausdrücklich vertagt. Sie ist im Buch nicht wiedergekommen.

Was hier zusammenkommt, ist die philosophische Form der physikalischen Asymmetrie. Hegel hat die Zeit das angeschaute Werden genannt, das Umschlagen von Sein in Nichts und Nichts in Sein; Bonsiepen hat gezeigt, dass Hegel den Raum dagegen substanzmetaphysisch, als ruhendes Sein denkt.\footnote{Bonsiepen, Hegels Raum-Zeit-Lehre, zu den Nachschriften und zur Genese.} Die Zeit gehört auf die Seite des Nichts, des Modus, der Subjektivität; der Raum auf die Seite des Seins, des Gegenstandes, der Sache. Die Physik, die vom Nichts absieht, muss die Zeit deshalb nach dem Raum behandeln, und sie tut es mit Recht, denn sie will die Sache. Was sie dabei verliert, ist nicht ein Teil der Sache, sondern der Modus, in dem die Sache je gegenwärtig ist. Und dieser Modus ist, was die Zeit von allem unterscheidet, das man nebeneinanderlegen kann.

\chapter{Das dauernde Ich}

\section{Drei Zeiten}

Baumann hat, um die Zeit \zitat{nach ihrer eigenen Art} zu behandeln, drei Schichten unterschieden. Die erste ist die psychologische Zeit. Das bloße Nacheinander von Vorstellungen wäre noch keine Zeitvorstellung. \zitat{Es gehört zur Zeitvorstellung ausser dem Nacheinander der Vorstellungen etwas, das sich dieses Nacheinanders als solchen bewusst wird [\ldots], etwas Bleibendes in der Aufeinanderfolge der Ideen. Dies Bleibende ist in uns unsere Ichvorstellung.}\footnote{Baumann, Lehren, Bd.\,2, S.\,659.} Die Zeit entsteht, indem das Nacheinander auf das dauernde Ich bezogen wird, und die Dauer des Ich ist selbst nicht Zeit; sie liegt der Zeit voraus, näher der Ewigkeit als der Zeit. Kants Satz wird umgedreht: \zitat{das ›Ich denke‹ begleitet nicht alle meine Vorstellungen, sondern alle meine Vorstellungen schliessen sich an das Ich denke an.}\footnote{Baumann, Lehren, Bd.\,2, S.\,662.}

Die zweite Zeit ist die gewöhnliche, psychologisch-astronomische. Sie ist uns gegeben an unserem Leibe und seiner Wechselwirkung mit den äußeren Dingen: Tag und Nacht, Wachen und Schlaf, die regelmäßigen Ereignisse der Natur. Die dritte ist die astronomische Zeit, die Sternzeit, die ihre Gleichförmigkeit aus der Annahme unveränderter Erdrotation gewinnt; aus ihr entsteht \zitat{jenes Idealbild der Zeit}, das gemeinhin als die Zeit schlechtweg gilt. \zitat{Dazu kam, dass man diese astronomische Zeit als eine Parallele zum Raum ansah, während beide Vorstellungen gar nicht mit einander stehen und fallen.}\footnote{Baumann, Lehren, Bd.\,2, S.\,663--667.} Kants Zeitlehre bedarf deshalb der Berichtigung: Er ging vom astronomischen Idealbild aus und wollte es als reine Anschauung gleich der des Raumes erweisen. Und die Gleichstellung von Raum und Zeit ist \zitat{durchaus unzulässig}.\footnote{Baumann, Lehren, Bd.\,2, S.\,668.}

Man kann Baumanns drei Zeiten als eine Ordnung der Parallelisierung lesen. Die dritte Zeit, das Idealbild, ist die Zeit als Linie; sie ist es, die man nach dem Raum behandelt hat. Die zweite ist die Zeit am Leib, die Zeit, die an Tag und Nacht abgelesen wird; ihr entspricht in der Physik die Zeit längs der Weltlinie eines Dinges, die Reichenbach aus der Genidentität gewinnt. Die erste ist die Zeit, die auf das Ich bezogen ist, dessen Dauer nicht Zeit ist. Je weiter man von der Linie zurückgeht, desto weniger gleicht die Zeit dem Raum.

\section{Das ewige Jetzt}

Reichenbach hat am Anfang des Zeit-Abschnitts einen Satz geschrieben, den man in einem Buch über die Physik der Raum-Zeit nicht erwartet: \zitat{›ich bin‹ ist immer gleichbedeutend mit ›ich bin jetzt‹; aber ich bin in einem ›ewigen Jetzt‹}, und fühle mich identisch bleibend.\footnote{Reichenbach, Raum-Zeit-Lehre, \S\,16.} Das Jetzt, in dem ich bin, ist immer dasselbe Jetzt, obwohl es immer ein anderes ist. Reichenbach vertagt die Erörterung dieser Ich-Zeit und kommt nicht auf sie zurück. Aber der Satz ist genau Baumanns Satz: Das Ich dauert nicht in der Zeit, sondern die Zeit wird auf das dauernde Ich bezogen. Und er steht nahe bei Brentanos Satz vom Beharren als sich erneuerndem Jetzt.

Hegel hat es in der Enzyklopädie gesagt, und Koch hat die Stelle hervorgehoben: Die Zeit ist dasselbe Prinzip wie das Ich = Ich des reinen Selbstbewusstseins, aber dasselbe in seiner Äußerlichkeit.\footnote{Hegel, Enzyklopädie \S\,258, zitiert nach Koch, Hegelsche Subjekte in Raum und Zeit.} Die Zeit ist das Ich, das sich nicht bei sich hat; das Ich ist die Zeit, die sich bei sich hat. Drei Autoren, die sich an dieser Stelle nicht aufeinander berufen, ein Physiker, ein Realist, ein Idealist, sagen dasselbe: Dass das Beharren, das die Zeit möglich macht, nicht im Raum zu suchen ist, sondern in der Dauer dessen, für den etwas jetzt ist.

\section{Das Beharrliche im Raum}

Damit ist nun Kants Umkehrung zu bedenken, die in der Widerlegung des Idealismus steht: Die Bestimmung meiner Existenz in der Zeit setzt ein Beharrliches voraus, das nicht in mir ist. Das Beharrliche ist die Substanz, und die Substanz ist im Raum. Martić hat gezeigt, wie ernst Kant das nimmt: Ohne Raum würde die Zeit selbst nicht als Größe vorgestellt werden; die Zeit wird nur am Beharrlichen zur messbaren Reihe.\footnote{Martić, Zeit und Raum, Kapitel zur Widerlegung des Idealismus und zu den Analogien.}

Die beiden Sätze, Baumanns und Kants, scheinen sich zu widersprechen: Das Beharren liegt im Ich; das Beharrliche liegt im Raum. Sie widersprechen sich nicht, wenn man zwei Fragen auseinanderhält. Kant fragt, woran die Zeit \emph{bestimmt}, das heißt gemessen und geordnet wird. Baumann fragt, woher die Zeit ihr \emph{Beharren} hat, das heißt: was es ist, das da bleibt, während die Vorstellungen wechseln. Auf die erste Frage ist die Antwort das Beharrliche im Raum, der Körper, an dem die Uhr abgelesen wird; das ist Reichenbachs Messkörper und Baumanns zweite Zeit. Auf die zweite Frage ist die Antwort das dauernde Ich. Das Beharrliche im Raum ist die äußere Gestalt der Dauer: Es ist, was die Dauer wird, wenn sie als Gegenstand angesehen wird. Brentano hat den Übergang beschrieben: Ein Ding, das besteht, beharrt durch ein sich wiederholendes Sich-Erneuern, und das ist die Weise, wie das Ich dauert, auf ein Räumliches übertragen.

Man muss hier vorsichtig sein. Es ist damit nicht gesagt, dass der Raum aus dem Ich stamme, oder das Beharrliche aus der Dauer. Es ist gesagt, dass das, was am Beharrlichen das Beharren ist, seine Weise zu sein von der Zeit hat und die Zeit die ihre vom dauernden Ich. Was der Raum ist, in dem das Beharrliche ist, bleibt dabei ungesagt. Reichenbach brauchte ihn als Messkörper, Kant als Substanz, Baumann als \zitat{Wesen sui generis}, das nicht in die Kategorientafel passt und das uns bleibt, auch wenn wir die Dinge wegdenken.\footnote{Baumann, Lehren, Bd.\,2, S.\,655; zum Raum als Anschauung, die ohne Dinge bleibt, S.\,648.} Der Raum widersteht der Ableitung an genau dieser Stelle, und die Stelle wird am Ende noch einmal aufgesucht.

\chapter{Der Übergang und seine Richtung}

\section{Zwei entgegengesetzte Übergänge}

Zwei der stärksten Gedanken, die von Zeit und Raum handeln, behaupten Übergänge in entgegengesetzter Richtung. Hegel: Der Raum geht in die Zeit über. Die Wahrheit des Raumes ist die Zeit; der Punkt, der im Raum keine Wirklichkeit hat, hat sie in der Zeit; \zitat{die Zeit ist also erst der wahrhafte Punkt}. Guzzoni, mit Heidegger, Wagner und Hölderlin: Die Zeit wird zum Raum. Sie eilt hin zum Ort, sie weilt und gewinnt Fläche, sie ist immer schon eingeräumt. Wer beide Sätze nebeneinanderstellt, sieht einen Widerspruch. Wer fragt, wovon sie jeweils sprechen, sieht keinen.

Hegels Übergang ist ein begrifflicher. Das Abstraktere hat seine Wahrheit im Konkreteren. Der Raum, das gleichgültige Nebeneinander, ist abstrakter als die Zeit, das Werden; deshalb ist die Zeit seine Wahrheit. Der Übergang sagt nichts darüber, was in der Erfahrung zuerst da ist oder woraus etwas entsteht; er ist kein Vorgang. Bonsiepen hat gezeigt, dass Hegel seit der Dissertation von 1801 dem abstrakten Raum die Subjektivität hinzufügen muss, den Punkt, die Zeit, und dass der Punkt, im Jenaer Systementwurf, \zitat{mens} heißt: Der Übergang vom Raum in die Zeit ist der Übergang von der Sache zum Geist.\footnote{Bonsiepen, Hegels Raum-Zeit-Lehre, zur Genese.}

Guzzonis Übergang ist einer der Erfahrung. Wenn die Zeit weilt, wird sie Raum: Sie breitet sich aus, sie hat Gegend, sie lässt sich aufhalten. Der Ort ist der Augenblick, an dem die Zeit sich niedergelassen hat. Dieser Übergang ist ein Vorgang, aber kein zeitlicher Vorgang in der Zeit; er ist die Weise, wie die Zeit ist, wenn sie nicht rennt.

\begin{figure}[htbp]
\centering
\begin{tikzpicture}[scale=1,>=Stealth,every node/.style={font=\small}]
% obere Reihe: Hegels räumlicher Aufstieg
\node (p) at (0,2) {Punkt};
\node (l) at (2.6,2) {Linie};
\node (f) at (5.2,2) {Fläche};
\node (r) at (7.8,2) {Raum};
\draw[->] (p) -- (l); \draw[->] (l) -- (f); \draw[->] (f) -- (r);
\node[gray] at (3.9,2.55) {\scriptsize der Punkt hat im Raum keine Wirklichkeit};
% Übergang Raum -> Zeit
\node (z) at (7.8,0) {Zeit};
\draw[->,thick] (r) -- node[right] {\scriptsize Wahrheit des Raumes} (z);
% untere Reihe: Zeit -> Jetzt -> Ort
\node (j) at (5.2,0) {Jetzt};
\node (o) at (2.6,0) {Ort};
\node (g) at (0,0) {Gegend};
\draw[->] (z) -- node[above] {\scriptsize \zitat{der wahrhafte Punkt}} (j);
\draw[->] (j) -- node[above] {\scriptsize räumliches Jetzt} (o);
\draw[->] (o) -- node[above] {\scriptsize Weile} (g);
% Rückkehr: Ort ist im Raum
\draw[->,thick,dashed] (o) -- node[left] {\scriptsize \zitat{immer schon eingeräumt}} (p);
% Beschriftung der Richtungen
\node[gray,align=left] at (10.4,2) {\scriptsize Hegel:\\[-2pt]\scriptsize der Raum geht\\[-2pt]\scriptsize in die Zeit über};
\node[gray,align=left] at (10.4,0) {\scriptsize Guzzoni:\\[-2pt]\scriptsize die Zeit wird\\[-2pt]\scriptsize zum Raum};
\end{tikzpicture}
\caption{Der Übergang und seine Richtung. Oben der begriffliche Aufstieg, in dem der Punkt sich im Raum ausbreitet und der Raum in die Zeit übergeht (Hegel). Unten der Weg der Erfahrung, in dem die Zeit als Jetzt zum Ort und als Weile zur Gegend wird (Guzzoni). Die gestrichelte Rückkehr: Der Ort ist ein Punkt im Raum, die Zeit ist eingeräumt.}
\label{abb:uebergang}
\end{figure}

\section{Dieselbe Rangordnung}

Beide Übergänge sagen dasselbe über den Rang. Bei Hegel ist die Zeit das Wahrhafte, der Raum die \zitat{leere Zerstreuung des Außersichseins}. Bei Guzzoni ist die Zeit \zitat{das Vorrangige und Bestimmende}, der Raum das, worin sie eingeräumt ist. Bei Reichenbach ist die Zeit das logisch Primäre, der Raum das Nebeneinander der Fasern. Bei Brentano ist der Begriff des Zeitlichen dem des Räumlichen übergeordnet. Bei Chiba ist die Zeit die weitere Klasse. Bei Baumann ist die Dauer des Ich das, ohne welches die Zeitvorstellung nicht gebildet werden könnte, und der Raum eine Anschauung, die uns ohne Dinge bleibt, aber nichts hervorbringt.

Die Übereinstimmung geht weiter als der Rang. Sie betrifft die \emph{Richtung der Ableitbarkeit}. Aus der Zeit erhält man den Raum: Die Gleichzeitigkeit gibt das Nebeneinander, das Jetzt gibt den Ort, die Lichtlaufzeit gibt die Länge, das Nahwirkungsprinzip gibt die Dimensionszahl. Aus dem Raum erhält man keine Zeit. Zenon scheitert daran, den Übergang aus Orten zusammenzusetzen. Brentanos Zeit mit einem einzigen Präteritalmodus wäre ein Topoid, also Raum. Und Leibniz' Definition des Raumes als Ordnung des Zugleichseienden ist, wie Baumann gezeigt hat, zirkulär: Das Zugleich ist zeitlich, und um daraus den Raum zu gewinnen, muss man \zitat{den Raum bereits mitbringen}; \zitat{Im Grunde lautet die Definition: Raum ist die Ordnung des Räumlichen.}\footnote{Baumann, Lehren, Bd.\,2, Kapitel zu Leibniz' Raum- und Zeitlehre.} Man kann also den Raum durch die Zeit definieren, und die Definition ist leer, weil sie das Räumliche voraussetzt; und man kann die Zeit nicht durch den Raum definieren, weil die Definition ein Topoid ergibt. Das ist die Asymmetrie in ihrer logischen Form.

\section{Der Zeitraum}

Kant hat in der Analytik beides gesehen. Martić fasst zusammen: Der Raum setzt die Zeit voraus und ist nur als Zeitraum konzipierbar; die Zeit wird nur am Raum zur Größe. Er zitiert Schelling: \zitat{Die Zeit wird nur durch den Raum, der Raum nur durch die Zeit endlich.}\footnote{Martić, Zeit und Raum, Kapitel zum Schematismus und zur Rezeption bei Fichte und Schelling.} Das Wechselverhältnis ist also nicht Parallelisierung. Es ist eine Ordnung mit einer Richtung und einem Vorbehalt. Die Richtung: Die Zeit stiftet die Ordnung, der Raum ist Zeitraum. Der Vorbehalt: Ohne Raum keine Größe, ohne Beharrliches keine Bestimmung, ohne Messkörper keine Lichtgeometrie.

Dass die beiden Übergänge sich zu diesem Wechselverhältnis fügen, ist das Ergebnis. Hegels Aufstieg vom Raum zur Zeit sagt, was das Wahrhafte ist; Guzzonis Weg von der Zeit zum Ort sagt, wo es sich aufhält. Beide zusammen sagen: Die Zeit ist das, was der Raum in Wahrheit ist, und der Raum ist das, worin die Zeit sich aufhält. Keine der beiden Seiten kann ohne die andere sein. Aber sie sind nicht gleich. Die eine gibt, die andere nimmt auf.

\chapter{Leib und Bewegung}

\section{Räumen und Zeitigen}

Fink hat den Vorrang, um den in der Tradition gestritten wird, zurückgewiesen, und zwar nach beiden Seiten. Kein Vorrang des Raumes, wie ihn die aristotelische Fundierung nahelegt, in der die Zeit als Zahl der Bewegung auf die Größe zurückgeht. Kein Vorrang der Zeit, wie ihn die Seelenlehre und dann Heideggers Fundamentalontologie behaupten. Stattdessen ein Vorrang der Welt, die als Zeit-Raum waltet, und ihr Walten ist die Bewegung: \zitat{Gesetzt den Fall, Bewegung sei ursprünglich das Räumen und Zeitigen}, dann müsste alles Bewegte in beidem sein, ohne dass eines das andere trüge.\footnote{Fink, Frühgeschichte, Kapitel zur Weltbewegung.} Die Bewegung ist bei Fink nicht ein Vorgang in Raum und Zeit, sondern das, was Raum und Zeit erst auseinanderlegt: Räumen wäre dann das Einräumen einer Weite, Zeitigen das Zumessen einer Weile.

Das ist ein Gedanke gegen alles, was bisher über den Rang der Zeit gesagt wurde, und er soll nicht verkleinert werden. Aber Fink selbst hält den Vorrang der Welt nicht ohne einen leisen Restvorrang der Zeit durch. Bei Platon liest er die Chora als das Raumhafte, die dunkle Materie, die den Ideen Platz gewährt, und den Nous als die Lichtmacht, die er in die Nähe der Zeit rückt; die Idee ist das Tun, die Chora das Leiden.\footnote{Fink, Frühgeschichte, Kapitel zu Platon.} Und die Bewegung, die er als Räumen und Zeitigen ansetzt, beschreibt er selbst so: \zitat{Bewegung ist die seltsame Weise, wo ein Ding im selben Moment an zwei Orten ist, d.\,h. im Übergang zwischen zwei Orten sich befindet.} Der Moment ist eines, die Orte sind zwei. Das Bewegte hält sich in einem Jetzt an zwei Orten auf, nicht an einem Ort zu zwei Zeiten. Das Räumen geschieht im Zeitigen, nicht das Zeitigen im Räumen. Finks Gleichrangigkeit hat also ihre eigene Richtung, und sie ist dieselbe wie die der anderen.

\section{Das Bewegte als Einheit}

Was Fink Welt nennt, nennt Hegel Materie und Bewegung. In der Nachschrift wird die Mechanik in zwei Teile geteilt: Raum und Zeit, und ihre Beziehung; die Beziehung hat zwei Gestalten, die Materie als relative Einheit und die Bewegung als absolute Einheit von Raum und Zeit.\footnote{Hegel, Naturphilosophie, Nachschrift, nach Bonsiepen.} Raum und Zeit können die Dialektik von Diskretheit und Kontinuität nicht an ihnen selbst verwirklichen; sie brauchen das, was in ihnen ist. Die Einheit von Zeit und Raum ist also nicht ein Drittes über ihnen, kein Behälter, in dem beide lägen, sondern das Bewegte, in dem beide zugleich sind.

Reichenbach sagt es mit der Genidentität: Längs der Zeit liegen die Zustände desselben Dinges, längs des Raumes verschiedene Dinge. Das Ding ist das, was die Faser bildet, und die Faser ist das, was die Zeit von der Zeit her und den Raum vom Nebeneinander her bestimmt. Ohne das Ding gäbe es weder Faser noch Nebeneinander, sondern nur einen Parameterraum, in dem nichts wirkt. Und die Parallelisierung von Raum und Zeit, die auf der Ebene der Mannigfaltigkeit so vollkommen erscheint, ist in Wahrheit die Abstraktion von dem, was in der Mannigfaltigkeit bewegt ist.

\section{Der Leib}

Unter den Bewegten ist eines ausgezeichnet, der Leib, und zwar nicht, weil er ein besonderer Körper wäre, sondern weil er das Bewegte ist, das sich selbst bewegt und das dabei jetzt und hier sagt. Koch hat gezeigt, dass er dem Raum die Richtungen gibt, Oben und Unten im Fallen, Vorn und Hinten im Gehen, Rechts und Links im Greifen. Baumann hat gezeigt, dass er die zweite Zeit gibt, die Zeit an Tag und Nacht, Wachen und Schlaf. Brentano hat gezeigt, dass er die dreidimensionale Grenze eines zeitlichen Kontinuums ist, das sich in ihm erneuert. Und Koch hat gezeigt, dass er das Partikulare ist, das sich a priori selbst individuiert und lokalisiert und damit die Spezifizierbarkeit aller anderen garantiert. Der Leib ist die Stelle, an der die Zeit im Raum ist: das Hier, das kein Dort ist, das Jetzt, das keine Konvention ist.

Das Verhältnis, das Koch als Wechselverhältnis von Subjektivität und raumzeitlicher Einzelnheit beschreibt, keine Seite ohne die andere, hat hier seinen Ort. Die Subjektivität braucht die raumzeitliche Einzelnheit, weil sie sonst kein Hier und kein Jetzt hätte, an dem sie sich individuiert. Die raumzeitliche Einzelnheit braucht die Subjektivität, weil sie sonst unspezifizierbar wäre, ein Universum von Doppelgängern. Und dieses Wechselverhältnis ist nicht symmetrisch, so wenig wie das von Zeit und Raum: Die Subjektivität ist der Ursprung des Werdens, die Einzelnheit sein Ergebnis.

\section{Die Weile im Fluss}

Guzzoni hat der Bewegung eine Gestalt gegenübergestellt, die die Tradition übersehen hat: das Bleiben. Die aristotelische Definition der Zeit als Zahl der Veränderung nach dem Früher und Später ist einseitig, denn nicht nur die Veränderung, auch die Ruhe hat Zeitcharakter. Die Weile ruht gleichsam im Fluss. Sie verläuft nicht linear, sie hat Anfang und Ende als eines zusammengesehen, sie ist offene Geschlossenheit. Und die Ruhe ist nicht das Fehlen von Bewegung, sondern ein \zitat{Raum für Bewegung}.\footnote{Guzzoni, Weile und Weite, Kapitel zur Weile und zur Ruhe.}

\begin{figure}[htbp]
\centering
\begin{tikzpicture}[scale=1,>=Stealth]
% Fluss der Zeit
\draw[->,very thick,decorate,decoration={snake,amplitude=1.2pt,segment length=9pt}] (-0.5,0) -- (10.5,0);
\node[below] at (10.2,-0.15) {\small Fluss der Zeit};
% Jetztpunkte
\foreach \x in {0.5,1.3,2.1,2.9} {\fill (\x,0) circle (0.05);}
\node[below] at (1.7,-0.2) {\scriptsize Jetzt, Jetzt, Jetzt \ldots};
% Weile als Fläche
\fill[gray!25] (4.2,-1.1) rectangle (7.6,1.1);
\draw[thick] (4.2,-1.1) rectangle (7.6,1.1);
\node at (5.9,0.55) {\small Weile};
\node at (5.9,-0.55) {\scriptsize keine halbe, keine doppelte};
% Ort als Augenblick des Raumes
\fill (5.9,0) circle (0.09);
\node[above right] at (5.95,0.05) {\scriptsize Ort};
\draw[<->,thin] (4.3,-1.35) -- (7.5,-1.35);
\node[below] at (5.9,-1.4) {\scriptsize Gegend};
% weiter Fluss
\foreach \x in {8.4,9.2} {\fill (\x,0) circle (0.05);}
\node[above] at (5.9,1.2) {\small \zitat{Diesem Fluß der Zeit gegenüber erscheint die Weile eher als eine Fläche oder ein Raum}};
\end{tikzpicture}
\caption{Die Weile im Fluss. Gegen die Folge der Jetztpunkte hat die Weile Fläche; in ihr ist der Ort der Augenblick des Raumes und die Gegend seine Weile (nach Guzzoni). Das Bleiben ist nicht Abwesenheit von Bewegung, sondern das, worin sie Raum hat.}
\label{abb:weile}
\end{figure}

Wenn Fink recht hat, dass die Bewegung Räumen und Zeitigen ist, dann hat Guzzoni recht, dass das Bleiben Einräumen und Weilen ist. Beide sind das Bewegte, in dem Zeit und Raum zugleich sind; das eine im Übergang, das andere im Aufenthalt. Und in beiden zeigt sich dieselbe Richtung. Im Übergang ist das Ding in einem Moment an zwei Orten: Die Zeit trägt, der Raum wird durchlaufen. Im Aufenthalt weilt die Zeit und gewinnt Fläche: Die Zeit wird Raum, der Raum nimmt sie auf. Es ist immer die Zeit, die etwas tut, und der Raum, dem es widerfährt. Finks Formel für Platon, Idee als Tun und Chora als Leiden, sagt es für den Ursprung des Gedankens. Aber Chora bedeutet auch das, ohne das nichts geschehen könnte. Das Leiden ist kein Mangel.

\chapter{Zahl und Vorzeichen}

\section{Die Zahl zwischen beiden}

Aristoteles hat die Zeit die Zahl der Bewegung nach dem Früher und Später genannt. Fink liest die Definition als Fundierung, Größe, Bewegung, Zeit, und bemerkt, dass die Zahl die Bewegung \zitat{umfasst}, wie der Ort als Gefäß den Körper umfasst.\footnote{Fink, Frühgeschichte, Kapitel zur aristotelischen Zeit-Analytik.} Die Zeit ist also bei Aristoteles nicht selbst eine Größe, sondern das, womit eine Größe, die Bewegung, gezählt wird. Damit steht die Zahl zwischen Zeit und Raum: Sie zählt das Räumliche, die durchlaufene Größe, und sie ist das Zeitliche, das Früher und Später.

Hegel hat die Zahl zwischen das Sinnliche und den Gedanken gestellt, und die Arithmetik, die man der Zeit als ihre Wissenschaft zuordnen möchte, reduziert die Zeit auf das tote Eins. Das Eins ist das Jetzt, aus dem das Werden herausgenommen ist; eine Zeit aus Einsen ist keine Zeit.\footnote{Bonsiepen, Hegels Raum-Zeit-Lehre, zur Frage einer Wissenschaft der Zeit.}

Baumann hat die gewöhnliche Ableitung der Zahl aus der Zeit, die sich auf die aristotelische Definition beruft, dreifach bestritten: \zitat{Erstens ist die Zahl discret, die Zeit continuirlich; sodann ist das blosse Nacheinander noch keine Zahl, so wenig wie es an sich schon Zeit ist, aber was in beiden Fällen hinzukommen muss, ist jedesmal etwas Anderes [\ldots]; endlich ist der Begriff der Einheit denkbar auch ohne die Zeit.}\footnote{Baumann, Lehren, Bd.\,2, S.\,668.} Das Eins ist ein geistiger Begriff, anwendbar wo und wie wir wollen; das Zusammenfassen von Eins und Eins zu Zwei ist ein geistiges Tun, das auch Gleichzeitiges betreffen kann. Die Zahl ist also weder aus der Zeit noch aus dem Raum, sondern aus dem Geist; und wo sie auf Zeit und Raum angewandt wird, verhält sie sich zu beiden verschieden. Den Raum lässt sie bestehen, denn seine Teile sind nebeneinander und lassen sich zählen, ohne dass einer vergeht. Die Zeit tötet sie, denn ihre Teile sind nicht nebeneinander, und wer sie zählt, hat sie nebeneinandergelegt. Das ist der Punkt, an dem die Parallelisierung entsteht: nicht in der Zeit und nicht im Raum, sondern in der Zahl, die beide gleich behandelt, weil sie beiden gleich fremd ist.

\section{Koordinaten und Zuordnung}

Reichenbach hat die Zahl in der Physik an ihren Ort gewiesen. Die Koordinaten sind bloße Zahlen; was sie bedeuten, wird durch Zuordnungsdefinitionen festgelegt, die sagen, welcher Körper starr ist, welche Uhr gleichförmig geht, welches Signal die Gleichzeitigkeit definiert. Die Zahl ist nicht das Maß, sie ist die Stelle, an die das Maß geschrieben wird. Schneider hat aus der allgemeinen Relativitätstheorie dasselbe gelesen: Die Koordinaten sind Symbole, messbar ist das Linienelement, und nur durch die Materie kommt physikalische Bedeutung hinein.\footnote{Reichenbach, Raum-Zeit-Lehre, \S\,4 und \S\,35; Schneider, Raum-Zeit-Problem, Kapitel zur allgemeinen Relativitätstheorie.}

Das heißt: Die Gleichbehandlung von Zeit und Raum, die in der Koordinatenmannigfaltigkeit vorliegt, ist eine Gleichbehandlung von Zahlen. Vier Zahlen legen ein Ereignis fest, das ist alles, was die Dimensionszahl zunächst besagt. Was die Zahlen bedeuten, ob sie Längen sind oder Zeiten, wird nicht von der Mannigfaltigkeit gesagt, sondern von der Zuordnung, und die Zuordnung ist für die eine Koordinate eine andere als für die drei übrigen.

\section{Das Vorzeichen}

Wo die Zuordnung anders ist, muss die Mannigfaltigkeit es irgendwo anzeigen. Sie zeigt es im Vorzeichen. Das Linienelement der Raum-Zeit ist keine Summe von Quadraten, sondern eine Differenz: Das Zeitglied hat das entgegengesetzte Vorzeichen der Raumglieder. Minkowski hat das Vorzeichen versteckt, indem er die Zeitkoordinate mit der imaginären Einheit multiplizierte; dann sieht das Linienelement aus wie eine euklidische Länge, und die Mannigfaltigkeit wie ein Raum mit einer Dimension mehr. Schneider referiert den Grund: Minkowski führte die Zeit als imaginäre Koordinate ein, um Symmetrie in der quadratischen Form zu erreichen. Die Symmetrie ist eine Symmetrie der Form; die Sache lässt sie, wie sie ist.\footnote{Schneider, Raum-Zeit-Problem, Kapitel zur speziellen Relativitätstheorie.}

Reichenbach nimmt das Vorzeichen ernst. Es ist für ihn der mathematische Ausdruck der Sonderstellung der Zeit, und die Sonderstellung ist die zuvor beschriebene: Die Zeit ist die Faserrichtung, die Dimension der Wirkungslinien, und nur sie. Dass man dieses Vorzeichen vor das Zeitglied oder, mit umgekehrter Konvention, vor die drei Raumglieder schreibt, ist Wahl. Dass es eines gibt, das ein Glied von drei anderen trennt, ist keine Wahl.\footnote{Reichenbach, Raum-Zeit-Lehre, \S\,27, \S\,41 und \S\,43.}

Hier ist eine Unterscheidung nötig, die leicht verwischt wird. Die Signatur der Metrik ist nicht die Asymmetrie. Sie ist deren metrische Spur, das, was von der Asymmetrie übrigbleibt, wenn man Zeit und Raum als Koordinaten schreibt. Die Asymmetrie selbst liegt vor der Metrik: in der Ordnung, in der Grenze, in der Richtung, im Modus. Das Vorzeichen zeigt an, dass diese Asymmetrie da ist; es sagt nicht, worin sie besteht. Wer die Signatur liest, als sagte sie, die Zeit sei reell und der Raum imaginär oder umgekehrt, liest eine Konvention als Befund. Was sich lesen lässt, ist nur dies: Eine Koordinate ist anders als drei.

Brentano hat die Rede von der Zeit als weiterer Raumdimension eine \zitat{unschädliche Fiktion} genannt.\footnote{Brentano, Untersuchungen, S.\,209.} Das trifft die Lage genau. Die Fiktion ist unschädlich, solange man weiß, dass sie eine ist; sie ist in der Physik unentbehrlich, weil die Physik die Sache will und der Modus in der Sache nicht vorkommt. Sie wird schädlich in dem Augenblick, in dem man das Vorzeichen für eine Auskunft darüber hält, was Zeit und Raum sind. Es ist die Auskunft, dass sie nicht dasselbe sind, und sonst nichts.

\section{Was die Zahl nicht erreicht}

Was die Zahl an der Zeit nicht erreicht, hat Brentano gesagt: den Modus, das Vergangen- und Künftigsein, das keine Bestimmung des Gegenstandes ist. Was sie am Raum nicht erreicht, hat Guzzoni gesagt: die Weite, die keine Höhe, Breite oder Länge ist. In beiden Fällen bleibt ein Unmaß, das dem Maß vorausgeht. Aber die beiden Unmaße sind nicht gleich. Die Weite ist der Raum, wie er ist, bevor man ihn misst; sie ist Raum. Der Modus ist die Zeit, wie sie ist, bevor man sie zählt; und er ist nicht Zeit im Sinne der Linie, sondern das, was die Linie zur Zeit macht. Die Zahl lässt vom Raum die Weite zurück und von der Zeit die Zeit.

\chapter{Was offen bleibt}

Die drei Asymmetrien sind gezeigt. Die Zeit hat eine ausgezeichnete Grenze, das Jetzt, und keine beliebigen; der Raum hat beliebige Grenzen und keine ausgezeichnete, und wo er eine hat, den Ort, das Hier, ist sie Zeit im Raum oder Leib im Raum. Die Zeit hat eine Richtung, weil sie die Faser der Wirkungen ist; der Raum hat keine, und wo er Richtungen zeigt, kommen sie vom Leib, der sich in ihm bewegt, bis hinab zur Zahl seiner Dimensionen, die der Sinn eines Wirkungsgesetzes ist. Die Zeit hat ihre Ordnung an der Kausalkette; der Raum hat seine Ordnung von der Zeit und sein Maß von der Uhr, und was die Uhr braucht, ist ein Ding, kein Raummaß. Der gemeinsame Grund ist, dass die Differenzen des Raumes Differenzen des Gegenstandes sind und die der Zeit Differenzen des Anerkennens. Die Physik, die vom Gegenstand handelt, behandelt die Zeit deshalb nach dem Raum, mit Recht und mit einem Verlust, der im Vorzeichen angezeigt wird und in ihm nicht enthalten ist.

Die Rangfrage ist damit beantwortet, sofern sie Grenze, Richtung und Maß betrifft. Die Zeit gibt, der Raum nimmt auf. Aus der Zeit erhält man die Ordnung des Raumes, aus dem Raum keine Ordnung der Zeit. Wo der Raum wahrhaft ist, ist er Zeit; wo die Zeit sich aufhält, ist sie Raum.

Was nicht beantwortet ist, ist die Frage nach dem Raum selbst. Alle Bestimmungen, die er bisher erhalten hat, sind negative oder geliehene: das Nebeneinander der Fasern, die leere Zerstreuung des Außersichseins, das Leiden, die Weite, das Eingeräumtsein, die Isotropie. Aber an drei Stellen hat er widerstanden, und diese Stellen sind nicht zufällig dieselben, an denen die Ableitung am weitesten getrieben wurde. Reichenbach, der die Raummessung auf die Zeitmessung zurückführt, braucht einen materiellen Messkörper, den die Lichtgeometrie nicht liefert. Kant, der die Zeit zur Form aller Erscheinungen macht, braucht ein Beharrliches im Raum, ohne das die Zeit nicht einmal als Größe vorgestellt werden könnte. Heidegger, der die Räumlichkeit auf die Zeitlichkeit zurückführen wollte, hat es zurückgenommen. Baumann, der die Zeit aus der Dauer des Ich versteht, lässt den Raum als eine Anschauung stehen, die uns bleibt, wenn wir die Dinge wegdenken, und die in keine Kategorientafel passt.

Und Fink, der die Bewegung als Räumen und Zeitigen zugleich denkt, lässt den Raum dem Ursprung nach gleichrangig.

Man kann diese Stellen so lesen, dass der Raum nur der Rest ist, der übrig bleibt, wenn die Ableitung noch nicht vollendet ist. Man kann sie auch so lesen, dass der Raum etwas ist, das aus der Zeit nicht abgeleitet werden kann, weil es nicht von ihrer Art ist: kein Modus, kein Werden, keine Faser, sondern das, worin Fasern nebeneinander sein können, ohne dass eine die andere wäre. Zwischen diesen beiden Lesarten wird hier nicht entschieden. Es wird nur festgehalten, dass die zweite nicht durch das widerlegt ist, was gezeigt wurde. Alles, was gezeigt wurde, betrifft die Grenze, die Richtung und das Maß des Raumes; es betrifft nicht das, was begrenzt, gerichtet und gemessen wird.

Offen bleibt auch das, was Reichenbach vertagt hat: die Einsinnigkeit der Zeit, das Fortschreiten, das \zitat{ein evidentes Erlebnis und zugleich eine Aussage über die Wirklichkeit} ist. Es ist nun zu sehen, warum die Physik es nicht lösen kann. Das Fortschreiten ist der Wechsel des Modus, das Künftige, das gegenwärtig wird und vergeht. Die Physik hat den Modus nicht, und sie kann ihn nicht haben, ohne aufzuhören, vom Gegenstand zu handeln. Dass die Zeit fortschreitet, ist deshalb keine physikalische Aussage und keine bloß subjektive; es ist die Aussage, dass es Gegenwart gibt, und Gegenwart ist das, was ein Gegenstand nicht ist und wovon jeder Gegenstand her ist. Chibas Satz sagt es einfach: Um die Zukunft zu erfahren, muss man nur warten. Warum das so ist, warum die Zukunft von selbst kommt, die Vergangenheit nie und das Dort nur, wenn man geht, ist mit den drei Asymmetrien nicht erklärt. Es ist das, was sie voraussetzen.

Und offen bleibt schließlich das Und. Fink hat gefragt, ob Zeit und Raum nebeneinander oder gleichzeitig sind. Die Antwort, die sich ergeben hat, lautet: weder noch. Sie sind zugleich in dem Bewegten, das räumt und zeitigt, und in dem Bleibenden, das einräumt und weilt. Das Bewegte und das Bleibende sind der Leib und die Dinge, und was diese sind, davon war hier nur insoweit die Rede, als sie Zeit und Raum auseinanderlegen. Die Welt, die Fink als das nennt, worin beide beisammen sind, ist mit ihren beiden Seiten nicht beschrieben. Sie ist das, worin die Beschreibung stattfindet.


% --- Apparat -----------------------------------------------------------
\appendix
\chapter{Statustafel}

Rang der Befunde. \emph{Belegt} heißt: an den herangezogenen Texten wörtlich ausgewiesen. \emph{Erschlossen} heißt: aus belegten Sätzen mehrerer Autoren zusammengeführt; die Zusammenführung ist eigene Deutung. \emph{Offen} heißt: nicht entschieden, absichtlich.

\begin{itemize}
\item Belegt: Das zeitliche Kontinuum besteht nur als Grenze, das räumliche als Ganzes (Brentano).
\item Belegt: In der Zeit ist ein Punkt der Sache nach ausgezeichnet, im Raum keiner (Brentano, S.\,117).
\item Belegt: Der Raum lässt sich überall begrenzen, aber nicht unterbrechen (Hegel, Nachschrift).
\item Belegt: Die Zeit ist der wahrhafte Punkt; der Ort ist das räumliche Jetzt (Hegel, Nachschrift und Enzyklopädie).
\item Belegt: Die Zeit ist die Faserrichtung der Kausalketten, der Raum das Nebeneinander der Faserbündel (Reichenbach, \S\,43).
\item Belegt: Die Zeit ist gegenüber dem Raum das logisch Primäre; das Lichtjahr ist die logische Urform aller Längenmessung (Reichenbach, \S\,27, \S\,42).
\item Belegt: Die Dimensionszahl des Raumes ist der objektive Sinn des Nahwirkungsprinzips (Reichenbach, \S\,44).
\item Belegt: Die Einsinnigkeit der Zeit bleibt bei Reichenbach ungelöst (\S\,21).
\item Belegt: Räumliche Differenzen sind sachlich, zeitliche modal (Brentano, VII--VIII); später abgeschwächt (Brentano, X).
\item Belegt: Das Jetzt ist weltweit, das Hier je meines (Fink).
\item Belegt: Wir selbst sind die notwendige Garantie der indexikalischen Spezifizierbarkeit (Koch 1994).
\item Belegt: Die Richtungen des Raumes kommen vom Leib, die Modi der Zeit von den Seelenvermögen (Koch 2010).
\item Belegt: Die Zeit wird zum Raum, nicht umgekehrt; die Zeit ist immer schon eingeräumt (Guzzoni, mit Heidegger).
\item Belegt: Leibniz' Definition des Raumes setzt das Räumliche voraus (Baumann).
\item Belegt: Die Zeitvorstellung wird auf das dauernde Ich bezogen, dessen Dauer der Zeit vorausgeht (Baumann, S.\,659--668).
\item Belegt: Die Zahl ist weder aus der Zeit noch aus dem Raum abzuleiten (Baumann, S.\,668--669).
\item Belegt: Der Raum setzt die Zeit voraus und ist nur als Zeitraum konzipierbar; die Zeit wird nur am Raum zur Größe (Martić, zu Kant).
\item Belegt: Um die Zukunft zu erfahren, muss man nur warten (Chiba).
\item Belegt: Das negative Vorzeichen ist Ausdruck der Sonderstellung der Zeit, nicht deren Aufhebung (Reichenbach; Schneider zu Minkowski).
\item Erschlossen: Der Chiasmus der Grenze, ein ausgezeichneter Schnitt ohne beliebige gegen beliebige Schnitte ohne ausgezeichneten; Fink und Baumann widersprechen einander darin nicht.
\item Erschlossen: Das weltweite Jetzt ist die Zeit in der Gestalt des Raumes und als solches definitorisch; das nicht-definitorische Jetzt ist das an der Faser.
\item Erschlossen: Die Richtungen des Raumes, bis zur Dimensionszahl, stammen aus der Wirkung, also aus dem Zeitlichen.
\item Erschlossen: Die Abstraktion vom Standpunkt ist die Abstraktion von der Zeit; der zeit-transzendente Standpunkt macht die Zeit zum Raum.
\item Erschlossen: Baumanns Ich-Dauer, Reichenbachs ewiges Jetzt und Hegels Ich = Ich sagen dasselbe.
\item Erschlossen: Hegels Übergang Raum--Zeit und Guzzonis Übergang Zeit--Raum widersprechen sich nicht; beide ergeben dieselbe Rangordnung.
\item Erschlossen: Die Signatur ist die metrische Spur der Asymmetrie, nicht die Asymmetrie.
\item Offen: Was der Raum ist, sofern er nicht Grenze, Richtung und Maß hat; ob er Rest der Ableitung oder eigener Art ist.
\item Offen: Die Einsinnigkeit der Zeit; warum die Zukunft von selbst kommt.
\item Offen: Die Bedeutung des Und; die Welt als das, worin beide beisammen sind.
\end{itemize}

\chapter{Abbildungen}

\begin{itemize}
\item Abbildung \ref{abb:grenze}: Ganzes und Grenze. Topisch und chronisch Kontinuierliches; der Chiasmus der Schnitte (nach Brentano).
\item Abbildung \ref{abb:faser}: Faser und Nebeneinander. Lichtkegel, Weltlinie als Faser, Faserbündel als Raum, Gleichzeitigkeitsflächen für verschiedene Werte von $\varepsilon$ (nach Reichenbach).
\item Abbildung \ref{abb:doppelgaenger}: Das symmetrische Universum. Doppelgänger, von außen ununterscheidbar, von innen indexikalisch unterschieden (nach Koch).
\item Abbildung \ref{abb:parameterraum}: Parameterraum und Koordinatenraum. Zwei Punkte auf einer Geraden oder ein Punkt in der Ebene; das Nahwirkungsprinzip zeichnet die Dimensionszahl aus (nach Reichenbach).
\item Abbildung \ref{abb:lichtjahr}: Das Lichtjahr als Urform der Länge. Signal, Reflexion, Rückkehr; Länge aus halber Laufzeit (nach Reichenbach).
\item Abbildung \ref{abb:uebergang}: Der Übergang und seine Richtung. Begrifflicher Aufstieg Punkt--Linie--Fläche--Raum--Zeit (Hegel) und Erfahrungsweg Zeit--Jetzt--Ort--Gegend (Guzzoni).
\item Abbildung \ref{abb:weile}: Die Weile im Fluss. Weile als Fläche gegen die Folge der Jetztpunkte; Ort als Augenblick des Raumes (nach Guzzoni).
\end{itemize}

\chapter{Sachen}

\begin{itemize}
\item Anerkennen, Modi des (Brentano): Kapitel 6.
\item Asymmetrie, drei Gestalten: Grenze (Kapitel 2), Richtung (Kapitel 4), Maß (Kapitel 5).
\item Beharren, Beharrliches: Kapitel 2, 7.
\item Bewegung als Räumen und Zeitigen: Kapitel 9.
\item Chora und Idee: Kapitel 9.
\item Dauer des Ich: Kapitel 7.
\item Dimensionszahl: Kapitel 4, 10.
\item Doppelgänger, symmetrisches Universum: Kapitel 3, 4.
\item Einsinnigkeit der Zeit: Kapitel 4, 11.
\item Faserrichtung, Faserbündel: Kapitel 3, 4.
\item Genidentität: Kapitel 4, 9.
\item Gleichzeitigkeit, $\varepsilon$-Definition: Kapitel 3.
\item Grenze, in recto und in obliquo: Kapitel 2.
\item Hier und Jetzt: Kapitel 3.
\item Idealbild der Zeit: Kapitel 7.
\item Indexikalität, Spezifizierbarkeit: Kapitel 3, 9.
\item Inkongruente Gegenstücke: Kapitel 4.
\item Kennzeichenprinzip: Kapitel 4.
\item Leib: Kapitel 2, 4, 9.
\item Lichtgeometrie, Lichtjahr: Kapitel 5.
\item Messkörper: Kapitel 5, 7, 11.
\item Nahwirkungsprinzip: Kapitel 4, 5.
\item Ort als räumliches Jetzt: Kapitel 2, 8.
\item Ort als Augenblick des Raumes: Kapitel 5, 9.
\item Parallelisierung: Kapitel 0, 1, 10.
\item Parameterraum, Koordinatenraum: Kapitel 4.
\item Punkt, wahrhafter: Kapitel 2, 8.
\item Sachliche und modale Differenzen: Kapitel 6.
\item Signatur, Vorzeichen: Kapitel 1, 10.
\item Topoid, Chronoid: Kapitel 6.
\item Topologie und Metrik: Kapitel 5.
\item Übergang, Raum in Zeit, Zeit in Raum: Kapitel 8.
\item Weile und Weite: Kapitel 5, 9, 10.
\item Zahl: Kapitel 10.
\item Zeit-transzendenter Standpunkt: Kapitel 6.
\item Zeitraum (Kant nach Martić): Kapitel 1, 8.
\item Zuordnungsdefinition: Kapitel 3, 5, 10.
\end{itemize}

\chapter{Personen}

\begin{itemize}
\item Aristoteles: Zeit als Zahl der Bewegung; Richtungen nur in beseelten Körpern; Ort als Gefäß. Kapitel 4, 5, 10.
\item Baumann, Julius: Kritik der Parallelisierung; drei Zeiten; Ich-Dauer; Zahl aus dem Geist. Kapitel 0, 2, 4, 7, 8, 10.
\item Bonsiepen, Wolfgang: Hegels Raum-Zeit-Lehre und ihre Genese; Edition der Nachschrift. Kapitel 1, 2, 5, 6, 8, 9.
\item Brentano, Franz: Kontinuum als Ganzes und als Grenze; sachliche und modale Differenzen; Überordnung des Zeitlichen. Kapitel 2, 5, 6, 7, 10.
\item Chiba, Kiyoshi: raumzeitliche Gegenstände; Warten und Gehen; zeit-transzendenter Standpunkt. Kapitel 6, 11.
\item Einstein, Albert: Relativitätstheorie, nach Schneider und Reichenbach. Kapitel 1, 5.
\item Fink, Eugen: das Und; Zenons Parallelisierung; Jetzt weltweit, Hier je meines; Räumen und Zeitigen. Kapitel 0, 1, 2, 3, 5, 9, 11.
\item Guzzoni, Ute: Weile und Weite; Ort als Augenblick des Raumes; die Zeit wird zum Raum. Kapitel 0, 2, 5, 8, 9, 10.
\item Hegel, Georg Wilhelm Friedrich: zwei Weisen des Außersichseins; die Zeit als wahrhafter Punkt; Ort als räumliches Jetzt; Zeit und Ich = Ich. Kapitel 0, 1, 2, 5, 6, 7, 8, 9.
\item Heidegger, Martin: Zeit-Raum; Rücknahme der Rückführung der Räumlichkeit auf die Zeitlichkeit, nach Guzzoni und Fink. Kapitel 0, 5, 9, 11.
\item Hölderlin, Friedrich: \zitat{Zeit eilt hin zum Ort}, nach Guzzoni. Kapitel 5.
\item Kant, Immanuel: Anschauungsformen; Schematismus; Widerlegung des Idealismus; inkongruente Gegenstücke; Dimensionszahl aus dem Kraftgesetz. Kapitel 1, 4, 7, 8, 11.
\item Koch, Anton Friedrich: Spezifizierbarkeit und Indexikalität; Richtungen aus dem Leib; Physik als Abstraktion von der Subjektivität. Kapitel 2, 3, 4, 6, 7, 9.
\item Leibniz, Gottfried Wilhelm: Raum als Ordnung der Koexistenzen, nach Baumann. Kapitel 8.
\item Martić, Marko: Kants synthetischer Vorrang der Zeit; Zeitraum; Newtons Anisotropie der Zeit. Kapitel 1, 4, 7, 8.
\item Minkowski, Hermann: Zeit als imaginäre Koordinate, nach Schneider und Reichenbach. Kapitel 1, 10.
\item Newton, Isaac: Anisotropie der Zeit, Isotropie des Raumes, nach Martić. Kapitel 4.
\item Platon: Chora und Idee, nach Fink. Kapitel 9.
\item Reichenbach, Hans: Zuordnungsdefinition; Kausaltheorie der Zeit; Faserrichtung; Lichtgeometrie; Dimensionszahl; ewiges Jetzt. Kapitel 1, 3, 4, 5, 6, 7, 8, 9, 10, 11.
\item Schelling, Friedrich Wilhelm Joseph: wechselseitige Endlichkeit von Zeit und Raum, nach Martić. Kapitel 8.
\item Schneider, Ilse: Raum und Zeit verlieren die physikalische Gegenständlichkeit; Koordinaten als Symbole. Kapitel 1, 4, 10.
\item Wagner, Richard: \zitat{Zum Raum wird hier die Zeit}, nach Guzzoni. Kapitel 5.
\item Zenon von Elea: Paradoxien als Parallelisierung, nach Fink. Kapitel 1, 8.
\end{itemize}

\chapter{Quellen}

\sloppy Herangezogene Texte, in der Reihenfolge ihrer Erscheinungsjahre.

\begin{itemize}
\item Baumann, Julius: Die Lehren von Raum, Zeit und Mathematik in der neueren Philosophie nach ihrem ganzen Einfluss dargestellt und beurtheilt. Bd.\,2: Leibniz, Leibniz und Clarke, Berkeley, Hume; Kurzer Lehrbegriff; Schluss und Regeln. Berlin 1869.
\item Schneider, Ilse: Das Raum-Zeit-Problem bei Kant und Einstein. Berlin 1921.
\item Reichenbach, Hans: Philosophie der Raum-Zeit-Lehre (1928). Gesammelte Werke, Bd.\,2, mit Erläuterungen von Andreas Kamlah. Braunschweig 1977.
\item Fink, Eugen: Zur ontologischen Frühgeschichte von Raum -- Zeit -- Bewegung. Den Haag 1957.
\item Brentano, Franz: Philosophische Untersuchungen zu Raum, Zeit und Kontinuum. Aus dem Nachlass herausgegeben von Stephan Körner und Roderick M. Chisholm, mit Anmerkungen von Alfred Kastil. Philosophische Bibliothek Bd.\,293. Hamburg 1976.
\item Bonsiepen, Wolfgang: Hegels Raum-Zeit-Lehre. Dargestellt anhand zweier Vorlesungsnachschriften. Mit Edition der Nachschrift. Hegel-Studien 20, 1985.
\item Koch, Anton Friedrich: Wozu noch Erste Philosophie? Zeitschrift für philosophische Forschung 48, 1994, Heft 4.
\item Koch, Anton Friedrich: Hegelsche Subjekte in Raum und Zeit. Hegel-Jahrbuch 2010.
\item Chiba, Kiyoshi: Kants Ontologie der raumzeitlichen Wirklichkeit. Versuch einer anti-realistischen Interpretation der Kritik der reinen Vernunft. Kantstudien-Ergänzungshefte. Berlin 2012.
\item Guzzoni, Ute: Weile und Weite. Zur nicht-metrischen Erfahrung von Zeit und Raum. Freiburg 2017.
\item Martić, Marko: Zeit und Raum. Eine Untersuchung zu Kants Auffassung der Anschauungsformen. Kantstudien-Ergänzungshefte 217. Berlin 2022.
\end{itemize}

Primärstellen, die über die genannten Texte herangezogen wurden: Aristoteles, Physik IV und De caelo II; Hegel, Enzyklopädie der philosophischen Wissenschaften (1830), \S\S\,254--261, sowie die Vorlesungsnachschriften nach Bonsiepen; Kant, Kritik der reinen Vernunft (B\,275, B\,291, B\,439), Von dem ersten Grunde des Unterschiedes der Gegenden im Raume (1768), Gedanken von der wahren Schätzung der lebendigen Kräfte (1747); Leibniz, Initia rerum mathematicarum metaphysica, nach Baumann; Heidegger, Sein und Zeit \S\,70, Zeit und Sein, Beiträge zur Philosophie, nach Guzzoni.

\chapter{Dokumentdaten}

\begin{itemize}
\item Titel: Zeit und Raum. Grenze, Richtung, Maß
\item Version: 2 (nach Review der Version 1)
\item Datum: 2026-09-25
\item Master-Datei: zeit-und-raum\_v2.tex
\item Kapitel-Dateien: kapitel/00\_das-und.tex bis kapitel/11\_was-offen-bleibt.tex
\item Apparat-Dateien: apparat/01\_statustafel.tex bis apparat/06\_dokumentdaten.tex
\item Konzept: konzept/konzept\_zeit-und-raum\_v2\_final.md (ersetzt v1)
\item Lesenotizen: lesenotizen/
\item Frühere Version: versionen/zeit-und-raum\_v1.tex, versionen/zeit-und-raum\_v1.pdf
\item Satz: pdflatex, scrbook, TikZ
\end{itemize}


\end{document}
